% concatenated from main-2.tex

\documentclass[11pt, reqno]{amsart}
\usepackage[utf8]{inputenc}
\usepackage{amsmath, amssymb, amsfonts, latexsym, amsthm}
\usepackage{calligra}
\usepackage{tcolorbox}
% \usepackage{refcheck}



\textwidth=455pt 
\evensidemargin=8pt \oddsidemargin=8pt
\marginparsep=8pt \marginparpush=8pt 
\textheight=640pt
\topmargin=-20pt

\newtheorem{thm}{Theorem}[section]
\newtheorem{cor}[thm]{Corollary}
\newtheorem{lem}[thm]{Lemma}
\newtheorem{prop}[thm]{Proposition}
\newtheorem{rem}[thm]{Remark}
\newtheorem{obs}[thm]{Observation}
\newtheorem{defn}[thm]{Definition}
\newtheorem{exam}[thm]{Example}
\numberwithin{equation}{section}

%\title{Weighted Szeg\H{o} kernels on planar domains}
%\author{}
%\date{Feb 2023}

\newcommand{\mbb}{\mathbb}
\newcommand{\ra}{\rightarrow}
\newcommand{\z}{\zeta}
\newcommand{\pa}{\partial}
\newcommand{\ov}{\overline}
\newcommand{\sm}{\setminus}
\newcommand{\ep}{\epsilon}
\newcommand{\no}{\noindent}
\newcommand{\al}{\alpha}
\newcommand{\Om}{\Omega}
\newcommand{\ti}{\tilde}
\newcommand{\la}{\lambda}
\newcommand{\ka}{\kappa}
\newcommand{\norm}[1]{\left\Vert#1\right\Vert}
\newcommand{\Norm}[1]{\left\Vert#1\right\Vert_{(1)}}
\newcommand{\Abs}[1]{\left\Vert#1\right\Vert_{(2)}}
\newcommand{\abs}[1]{\left\vert#1\right\vert}
\newcommand{\de}{\delta}
\newcommand{\De}{\Delta}
\newcommand{\La}{\Lambda}
\newcommand{\mf}{\mathbf}
\newcommand{\om}{\omega}
\newcommand{\ga}{\gamma}
\newcommand{\Ga}{\Gamma}
\newcommand{\be}{\beta}
\newcommand{\Si}{\Sigma}

\begin{document}
	\title[Weighted kernel functions on planar domains]{Weighted kernel functions on planar domains}
	\keywords{weighted Szeg\H{o} kernel, weighted Garabedian kernel, weighted Ahlfors map, weighted Carathéodory metric, Bergman kernel}
	\thanks{A part of the work was done when the first named author was an institute postdoctoral fellow at the Indian Institute of Technology Palakkad, India. The first named author was partially supported by SHENG-3 K/NCN/000286, research project UMO2023/Q/ST1/00048}
	\subjclass[2020]{Primary: 30C40; Secondary: 31A99}
	\author{Aakanksha Jain, Kaushal Verma} 
	\address{AJ: Department of Mathematics, Indian Institute of Technology, Kanjikode, Pudussery west, Palakkad, Kerala 678623, India}
    \address{Current address: Institute of Mathematics, Jagiellonian University, ul Lojasiewicza 6, 30-348 Kraków, Poland}
    
	\email{aakanksha.jain@uj.edu.pl, aakankshaj@alum.iisc.ac.in}
	
	\address{KV: Department of Mathematics, Indian Institute of Science, Bengaluru-560012, India}
	\email{kverma@iisc.ac.in}
	
\begin{abstract}
We study the variation of weighted Szeg\H{o} and Garabedian kernels on planar domains as a function of the weight. A Ramadanov type theorem is shown to hold as the weights vary. As a consequence, we derive properties of the zeros of the weighted Szeg\H{o} and Garabedian kernel for weights close to the constant function $1$ on the boundary. We further study the weighted Ahlfors map and strengthen results concerning its boundary behaviour. Explicit examples of the weighted kernels are presented for certain classes of weights. We highlight an interesting property of the weighted Szeg\H{o} and Garabedian kernels, implicit in Nehari’s work, and explore several of its consequences. Finally, we discuss the weighted Carathéodory metric, and describe relations of the weighted Szeg\H{o} and Garabedian kernel with certain classical kernel functions.
\end{abstract}
	
\maketitle

\section{Introduction}

\noindent Some aspects of weighted Szeg\H{o} kernels on a finitely connected planar domain were considered in \cite{JV}. We continue this study by strengthening some of the results obtained therein on the variation of the weighted Szeg\H{o} kernels as functions of the weight and noting some consequences about the zeros of the weighted Szeg\H{o} kernel. This has direct consequences about the zeros of the weighted Garabedian kernel. Continuing along this path leads to the consideration of a weighted Ahlfors map (cf. Nehari \cite{Nehari}) associated with the given planar domain, and we are able to strengthen an observation made by Nehari in \cite{Nehari} about its boundary behaviour. We then provide explicit formulae for the weighted Szeg\H{o} kernel, the weighted Garabedian kernel, and the weighted Ahlfors map for certain classes of weights. We conclude by highlighting a specific property about weighted Szeg\H{o} and Garabedian kernels that is of independent interest but is implicit in \cite{Nehari}, and exploring a few interesting consequences. Finally, we discuss the weighted Carathéodory metric, as well as expressions for the weighted Szeg\H{o} kernel and the weighted Garabedian kernel in terms of certain classical kernel functions.

\medskip

To make all this precise and fix the notation, let $\Omega \subset \mbb C$ be a smoothly bounded $n$-connected domain without isolated boundary points and $\varphi$ a positive real-valued $C^{\infty}$-smooth function on $\pa \Om$. Following \cite{Be0}, \cite{Be1}, let $L^2_{\varphi}(\Omega)$ be the Hilbert space of complex-valued functions on $\partial\Omega$ that are square integrable with respect to the arc length measure $\varphi\,ds$. Here, $ds$ is the differential of arc length and is given by  $ds = \vert z'(t) \vert \; dt$ where $z = z(t)$ is a smooth parametrization of $\partial \Omega$. In terms of the complex unit tangent vector function $T(z(t)) = z'(t) / \vert z'(t) \vert$, it is seen that $dz = z'(t) \; dt = T ds$. The inner product on $L^2_{\varphi}(\partial\Omega)$ is
\begin{equation}
\langle u,v\rangle_{\varphi}=\int_{\partial\Omega}u\bar{v}\,\varphi\,ds\quad \text{for }u,v\in L^2_{\varphi}(\partial\Omega).
\end{equation}
When $\varphi \equiv 1$, this reduces to the standard inner product $\langle u, v \rangle$. Note that $L^2_{\varphi}(\partial\Omega)=L^2(\partial\Omega)$ as sets. Also, $C^{\infty}(\partial\Omega)$ is dense in $L^2_{\varphi}(\partial\Omega)$ with respect to its norm. For $u\in C^{\infty}(\partial\Omega)$, the Cauchy transform of $u$ is 
\begin{equation}
(\mathcal{C}u)(z)=\frac{1}{2\pi i}\int_{\partial\Omega}\frac{u(\zeta)}{\zeta-z}d\zeta,\quad z\in\Omega
\end{equation}
which is holomorphic on $\Omega$. For $a \in \Omega$ and $z \in \partial \Omega$, let $C_a(z)$ be the conjugate of 
\[
\frac{1}{2 \pi i} \frac{T(z)}{z-a}.
\]
Then, $C_a$ is the Cauchy kernel which defines the Cauchy transform $\mathcal C$ in the sense that $\mathcal C u(z) = \langle u, C_z \rangle$. The weighted Cauchy kernel is defined to be $\varphi^{-1} C_a$ and the corresponding weighted Cauchy transform $\mathcal C_{\varphi}u$ satisfies
\[
(\mathcal C_{\varphi} u)(z) = \langle u, \varphi^{-1} C_z \rangle_{\varphi} = \langle u, C_z \rangle = (\mathcal C u)(z)
\]
and this shows that $\mathcal C_{\varphi} = \mathcal C$.

\medskip

Let $A^{\infty}(\Omega)$ denote the space of holomorphic functions on $\Omega$ that are in $C^{\infty}(\ov{\Omega})$. The Cauchy transform $\mathcal C$ maps $C^{\infty}(\partial\Omega)$ into $A^{\infty}(\Omega)$  and this allows the Cauchy transform to be viewed as an operator from $C^{\infty}(\partial\Omega)$ into itself. Let $A^{\infty}(\partial\Omega)$ denote the set of functions on $\partial\Omega$ that are the boundary values of functions in $A^{\infty}(\Omega)$. The Hardy space $H^2(\partial\Omega)$ is defined to be the closure in $L^2(\partial\Omega)$ of $A^{\infty}(\partial\Omega)$ and members of the Hardy space are in one-to-one correspondence with the space of holomorphic functions on $\Omega$ with $L^2$ boundary values. This Hardy space can be identified in a natural way with the subspace of $L^2_{\varphi}(\partial\Omega)$, equal to the closure of $A^{\infty}(\partial\Omega)$ in that space. Thus, there is no need to define $H^2_{\varphi}(\partial\Omega)$ separately. But we shall use the notation $H^2_{\varphi}(\partial\Omega)$ whenever there is a need to emphasize that the Hardy space is endowed with the inner product $\langle \cdot,\cdot\rangle_{\varphi}$. 

\medskip

The orthogonal projection $P_{\varphi}$ from $L^2_{\varphi}(\partial\Omega)$ onto $H^2(\partial\Omega)$ is the weighted Szeg\H{o} projection. For $a\in\Omega$, the evaluation functional $h\mapsto h(a)$ on $H^2_{\varphi}(\Omega)$ is continuous since
\[
\vert h(a)\vert = \vert \langle h, C_a \rangle \vert = \vert \langle h, C_a \varphi^{-1}\rangle_{\varphi}\vert \leq \Vert C_a \varphi^{-1} \Vert_{\varphi} \Vert h\Vert_{\varphi}
\]
and hence, there exists a unique function $S_{\varphi}(\cdot,a)\in H^2(\partial\Omega)$ such that 
\begin{equation}
h(a)=\langle h, S_{\varphi}(\cdot, a)\rangle_{\varphi}
\end{equation}
for all $h\in H^2(\partial\Omega)$. The function $S_{\varphi}(\cdot,\cdot)$ is the weighted Szeg\H{o} kernel of $\Omega$ with respect to the weight $\varphi$. 

\medskip

It can be seen that the weighted Szeg\H{o} kernel is hermitian symmetric, that is, for $z, w\in\Omega$, we have
$
S_{\varphi}(z,w)= \overline{S_{\varphi}(w,z)}.
$
Therefore, $S_{\varphi}(\cdot,\cdot)\in C^{\infty}(\Omega\times\Omega)$. 

\medskip

The weighted Szeg\H{o} kernel is the kernel of the weighted Szeg\H{o} projection because for all $u\in L^2_{\varphi}(\partial\Omega)$ and $a\in\Omega$
\begin{equation}
(P_{\varphi}u)(a) = \langle P_{\varphi}u, S_{\varphi}(\cdot, a)\rangle_{\varphi}=\langle u, S_{\varphi}(\cdot, a)\rangle_{\varphi} = \int_{\zeta\in\partial\Omega}S_{\varphi}(a,\zeta)u(\zeta)\varphi(\zeta) ds.
\end{equation}
Note that for all $h\in H^2(\partial\Omega)$
\[
h(a) = (\mathcal{C}h)(a) = \langle h, C_a\rangle = \langle h, \varphi^{-1}C_a\rangle_{\varphi} =  \langle h, P_{\varphi}(\varphi^{-1}C_a)\rangle_{\varphi}.
\]
Therefore, the weighted Szeg\H{o} kernel $S_{\varphi}(\cdot,a)$ is also given by the weighted Szeg\H{o} projection of $\varphi^{-1}C_a$.
That is,
\begin{equation}
S_{\varphi}(z,a)=(P_{\varphi} (\varphi^{-1}C_a))(z)=\int_{\zeta\in\partial\Omega}S_{\varphi}(z,\zeta)C_a(\zeta) ds
\end{equation}
for every $z\in\Omega$. In \cite{JV}, the following weighted Kerzman--Stein formula
\begin{equation}
P_{\varphi}(I+\mathcal{A}_{\varphi})=\mathcal{C}
\end{equation}
was shown to hold on $C^{\infty}(\partial\Omega)$, where $I$ is the identity and $\mathcal{A}_{\varphi} = \mathcal{C}-\mathcal{C}^{*}_{\varphi}$ is the weighted Kerzman-Stein operator. Since the Cauchy transform maps $C^{\infty}(\partial\Omega)$ into itself, it follows from the weighted Kerzman--Stein formula that the weighted Szeg\H{o} projection $P_{\varphi}$ also maps $C^{\infty}(\partial\Omega)$ into itself. Thus, $S_{\varphi}(\cdot,a)\in A^{\infty}(\Omega)$. In fact, $S_{\varphi}(z, w) \in \left( \overline \Om \times \overline \Om \right) \setminus \Delta$ where $\Delta = \{(z, z): z \in \pa \Omega\}$ is the diagonal in $\pa \Om \times \pa \Om$ as Theorem $3.2$ in \cite{JV} shows.

\medskip

A function $u$ in $L^2(\partial\Omega)$ is orthogonal to $H^2_{\varphi}(\partial\Omega)$ if and only if $u\varphi$ is orthogonal to $H^2(\partial\Omega)$. All the functions $v\in L^2(\partial\Omega)$ orthogonal to $H^2(\partial\Omega)$ are of the form $\overline{HT}$ where $H\in H^2(\partial\Omega)$. If $v\in C^{\infty}(\partial\Omega)$ then $H\in A^{\infty}(\Omega)$. So, the orthogonal decomposition of $\varphi^{-1}C_a$ is given by
\begin{equation}
\varphi^{-1}C_a = S_{\varphi}(\cdot,a) + \varphi^{-1} \overline{H_aT}
\end{equation}
where $H_a\in A^{\infty}(\Omega)$. Also, the above decomposition shows that $H_a$ is holomorphic in $a\in\Omega$ for fixed $z\in\partial\Omega$. The weighted Garabedian kernel of $\Omega$ with respect to $\varphi$ is defined by
\begin{equation}
L_{\varphi}(z,a) = \frac{1}{2\pi}\frac{1}{z-a} - i H_a(z).
\end{equation}
For a fixed $a\in\Omega$, $L_{\varphi}(z,a)$ is a holomorphic function of $z$ on $\Omega\setminus\{a\}$ with a simple pole at $z=a$ with residue $1/2\pi$, and extends $C^{\infty}$-smoothly to $\partial\Omega$. In addition, $L_{\varphi}(z,a)$ is holomorphic in $a$ on $\Omega$ for fixed $z\in \partial\Omega$. Moreover, it is known that (see \cite{Nehari})
\begin{equation}
L_{\varphi}(z,a) = - L_{1/\varphi}(a,z)
\end{equation}
for $z,a\in\Omega$. Therefore, for a fixed $z\in\Omega$, $L_{\varphi}(z,a)$ is a holomorphic function of $a$ on $\Omega\setminus\{z\}$ with a simple pole at $a=z$ and residue ${1}/2\pi$, and extends $C^{\infty}$-smoothly to $\partial\Omega$. Finally, for $z \in \partial \Omega$, $a \in \Omega$
\[
S_{\varphi}(a,z) = \overline{S_{\varphi}(a,z)} 
=
\frac{1}{\varphi(z)}\left(\frac{1}{2\pi i}\frac{T(z)}{z-a}-H_a(z) T(z)\right)
=
\frac{1}{i\varphi(z)}\left(\frac{1}{2\pi}\frac{1}{z-a}-iH_a(z)\right)T(z)
\]
shows that the weighted Szeg\H{o} kernel and the weighted Garabedian kernel satisfy the identity
\begin{equation}
S_{\varphi}(a,z) = \frac{1}{i\varphi(z)} L_{\varphi}(z,a) T(z).
\end{equation}
It follows from Theorem $3.3$ in \cite{JV} that the function $l_{\varphi}(z,w)$ defined by 
\[
L_{\varphi}(z,w) = \frac{1}{2\pi(z-w)} + l_{\varphi}(z,w)
\]
is holomorphic on $\Omega\times\Omega$ and extends to be in $C^{\infty} (\overline{\Omega}\times\overline{\Omega})$.


%%%%%%%%%%%%%%%%%%%%%%%%%%%%%%%%%%%%    
	\section{Convergence of weighted Szeg\H{o} kernels}

    The weights for which the weighted Szeg\H{o} kernel is well-defined are studied in \cite{pw-um}. We will focus only on a special class of weights that are smooth and positive functions on the boundary and study the continuity of the map $\varphi \mapsto S_{\varphi}$. In \cite{JV}, it was shown that for a sequence of positive real-valued $C^{\infty}$ functions $\varphi_k$ on $\partial\Omega$ such that $\varphi_k\rightarrow\varphi$ in the $C^{\infty}$-topology on $\partial\Omega$,  
		\begin{equation}
		\lim_{k\rightarrow\infty}S_{\varphi_k}(z,w) = S_{\varphi}(z,w)
		\quad
		\text{and}
		\quad
		\lim_{k\rightarrow\infty}l_{\varphi_k}(z,w) = l_{\varphi}(z,w)
		\end{equation}
		locally uniformly on $(\Omega\times\overline{\Omega})\cup(\overline{\Omega}\times\Omega)$. We extend this to the maximal possible set that is expected, in view of the fact that each Szeg\H{o} kernel is in $C^{\infty} (\overline{\Omega}\times\overline{\Omega}) \setminus \Delta$.   

	\begin{thm}\label{Szego}
		Let $\Omega\subset\mathbb{C}$ be a $C^{\infty}$-smoothly bounded $n$-connected domain and $\varphi$ a positive real-valued $C^{\infty}$ smooth function on $\partial\Omega$. 
		Let $\{\varphi_k\}_{k=1}^{\infty}$ be a sequence of positive real-valued $C^{\infty}$ functions on $\partial\Omega$ such that $\varphi_k\rightarrow\varphi$ in the $C^{\infty}$ topology on $\partial\Omega$ as $k\rightarrow\infty$. Then
		\begin{equation}
		\lim_{k\rightarrow\infty}S_{\varphi_k}(z,w) = S_{\varphi}(z,w) 
		\end{equation}
		locally uniformly on $(\overline{\Om}\times\overline{\Om})\setminus \Delta$. 
	\end{thm}
	
	\begin{proof}
		Let $z_0,w_0\in\partial\Omega$ be distinct arbitrary points. Choose relatively compact neighborhoods $U, V\subset\mathbb{C}$ of $z_0, w_0$ respectively such that they both intersect only one of the boundary components of $\Omega$ namely, the ones containing $z_0, w_0$, and satisfy $\overline{U}\cap \overline{V} = \emptyset$. Let $\Tilde{U} = U\cap\overline{\Omega}$ and $\Tilde{V} = V\cap\overline{\Omega}$. It is sufficient to show that 
		\[
		\lim_{k\rightarrow\infty}S_{\varphi_k}(z,w) = S_{\varphi}(z,w) 
		\]
		uniformly on $\Tilde{U}\times\Tilde{V}$. Let $f:\Omega\rightarrow \mathbb{D}$ be a proper holomorphic map which is an $n$-to-one branched covering and which maps each boundary curve bijectively to the unit circle. Suppose that $f$ has zeros at $a_1,\ldots,a_N \in \Omega$ with multiplicities $M(1),\ldots, M(N)$ respectively. Note that $M(1)+\cdots+M(N) = n$. For an arbitrary weight $\psi$, let $S_{\psi}^{(0)}(z, w)$ denote $S_{\psi}(z, w)$, and let $S_{\psi}^{\overline n}(z, w)$ denote $(\partial^n/\partial {\overline w}^n)S_{\psi}(z, w)$.
        By Bell \cite{Be1}, 
		\begin{equation}\label{rel_k}
			S_{\varphi_k}(z,w) = \frac{1}{1-f(z)\overline{f(w)}} \left(\sum_{i,j=1}^N \sum_{n=0}^{M(i)} \sum_{m=0}^{M(j)} c_{ijnmk}\, S_{\varphi_k}^{(\bar{n})}(z,a_i) \overline{S_{\varphi_k}^{(\bar{m})}(w,a_j)}\right) 
		\end{equation}
		where the constants $c_{ijnmk}$ are uniquely determined by the system of equations
		\begin{equation}\label{constants_k}
			\sum_{j=1}^N\sum_{m=0}^{M(j)}c_{ijnmk}\overline{S_{\varphi_k}^{(q,\bar{m})}(a_l,a_j)}
			=
			\left\{
			\begin{array}{lr}
				1 & \text{if } i=l \text{ and } n=q \\
				0 & \text{if } i\neq l \text{ or } n\neq q
			\end{array}
			\right.
		\end{equation}
		where $1\leq l\leq N$, $1\leq i\leq N$, $0\leq q\leq M(l)$ and $0\leq n\leq M(i)$.  Similarly, 
		\begin{equation}\label{rel}
			S_{\varphi}(z,w) = \frac{1}{1-f(z)\overline{f(w)}} \left(\sum_{i,j=1}^N \sum_{n=0}^{M(i)} \sum_{m=0}^{M(j)} c_{ijnm}\, S_{\varphi}^{(\bar{n})}(z,a_i) \overline{S_{\varphi}^{(\bar{m})}(w,a_j)}\right)
		\end{equation}
		where the constants $c_{ijnm}$ are uniquely determined by the system of equations
		\begin{equation}\label{constants}
			\sum_{j=1}^N\sum_{m=0}^{M(j)}c_{ijnm}\overline{S_{\varphi}^{(q,\bar{m})}(a_l,a_j)}
			=
			\left\{
			\begin{array}{lr}
				1 & \text{if } i=l \text{ and } n=q \\
				0 & \text{if } i\neq l \text{ or } n\neq q
			\end{array}
			\right.
		\end{equation}
		where $1\leq l\leq N$, $1\leq i\leq N$, $0\leq q\leq M(l)$ and $0\leq n\leq M(i)$.
		We can write the systems (\ref{constants_k}) and (\ref{constants}) of linear equations in matrix form as:
		\[
		\mathcal{A}_k\mathcal{C}_k = \mathcal{B}
		\quad {and} \quad
		\mathcal{A} \mathcal{C} = \mathcal{B}
		\]
		where $\mathcal{B}$ is an $n\times 1$ vector with entries $0$ or $1$, $\mathcal{A}_k$ is an $n\times n$ matrix with entries $0$ or $S_{\varphi_k}^{(q,\bar{m})}(a_l,a_j)$, 
		$\mathcal{A}$ is same as the matrix $\mathcal{A}_k$ if we replace the entries $S_{\varphi_k}$ by $S_{\varphi}$, and 
		$\mathcal{C}_k$ and $\mathcal{C}$ are $n\times 1$ vectors with entries $c_{ijnmk}$ and $c_{ijnm}$ respectively. The fact that $M(1) + \cdots +M(N) = n$ is being used here.
		
		\medskip
		
		Since $S_{\varphi_k}\rightarrow S_{\varphi}$ locally uniformly (and so its derivatives) on $\Omega\times\Omega$ as $k\rightarrow\infty$, we have that $\mathcal{A}_k^{-1}\rightarrow\mathcal{A}^{-1}$ uniformly as $k\rightarrow\infty$. Therefore, $c_{ijnmk}\rightarrow c_{ijnm}$ uniformly as $k\rightarrow\infty$.

        \medskip
		
		If $z_0$ and $w_0$ lie on the same boundary curve then $1-f(z)\overline{f(w)}$ does not vanish on $\Tilde{U}\times\Tilde{V}$ because $f$ is bijective on each boundary component of $\partial \Omega$. 

        \medskip
        
        When $z_0,w_0$ lie on different boundary curves say, $\gamma_1, \gamma_2$ respectively,  choose $z_1\in \gamma_1$ with $z_1\neq z_0$. By Bell--Kaleem \cite{BeKf}, there exists a proper holomorphic map $f: \Omega \rightarrow \mathbb D$ which maps both $z_1$ and $w_0$ to $1 \in \partial \mathbb D$ and is bijective on each boundary component of $\partial \Omega$. We can then shrink $U$, if necessary, to make sure that $1-f(z)\overline{f(w)}$ does not vanish on $\Tilde{U}\times\Tilde{V}$. 
        
		
		\medskip
		
		Since the zeros $a_i$ lie in $\Omega$, \cite{JV} implies that
		\[
		\lim_{k\rightarrow\infty}S_{\varphi_k}^{(\bar{n})}(z,a_i)= S_{\varphi}^{(\bar{n})}(z,a_i)
		\quad
		\text{and}
		\quad
		\lim_{k\rightarrow\infty}S_{\varphi_k}^{(\bar{m})}(w,a_j)= S_{\varphi}^{(\bar{m})}(w,a_j)
		\]
		uniformly for $z\in \tilde{U}$ and $w\in\tilde{V}$ respectively. These observations along with (\ref{rel_k}) and (\ref{rel}) show that 
		\[
		\lim_{k\rightarrow\infty}S_{\varphi_k}(z,w) = S_{\varphi}(z,w)
		\]
		uniformly on $\Tilde{U}\times \Tilde{V}$. The convergence of the derivatives also follows similarly.  
	\end{proof} 
	
	\begin{thm}
		Let $\Omega\subset\mathbb{C}$ be a $C^{\infty}$-smoothly bounded $n$-connected domain and $\varphi$ a positive real-valued $C^{\infty}$ smooth function on $\partial\Omega$. 
		Let $\{\varphi_k\}_{k=1}^{\infty}$ be a sequence of positive real-valued $C^{\infty}$ functions on $\partial\Omega$ such that $\varphi_k\rightarrow\varphi$ in $C^{\infty}$ topology on $\partial\Omega$ as $k\rightarrow\infty$. Then
		\begin{equation}
		\lim_{k\rightarrow\infty}l_{\varphi_k}(z,w) = l_{\varphi}(z,w)
		\end{equation}
		locally uniformly on $(\overline{\Om}\times\overline{\Om})\setminus\Delta$.
	\end{thm}
	
\begin{proof}
	As before, let $z_0,w_0\in\partial\Omega$ be distinct arbitrary points. Choose relatively compact neighborhoods $U, V\subset\mathbb{C}$ of $z_0, w_0$ respectively such that they both intersect only one of the boundary components of $\Omega$ namely, the ones containing $z_0, w_0$, and satisfy $\overline{U}\cap \overline{V} = \emptyset$. Let $\Tilde{U} = U\cap\overline{\Omega}$ and $\Tilde{V} = V\cap\overline{\Omega}$. It is sufficient to show that 
		\[
		\lim_{k\rightarrow\infty}l_{\varphi_k}(z,w) = l_{\varphi}(z,w) 
		\]
		uniformly on $\Tilde{U}\times\Tilde{V}$. Let $f:\Omega\rightarrow \mathbb{D}$ be a proper holomorphic map which is an $n$-to-one branched covering and which maps each boundary curve bijectively to the unit circle. Suppose that $f$ has zeros $a_1,\ldots,a_N$ with multiplicities $M(1),\ldots, M(N)$ respectively. 
        \medskip
        
        If $z_0$ and $w_0$ lie on the same boundary component of $\partial \Omega$, then $f(z)-f(w)$ does not vanish on $\Tilde{U}\times\Tilde{V}$ because $f$ is bijective on each boundary component of $\partial \Omega$. 

        \medskip
        
        When $z_0,w_0$ lie on different boundary curves, say, $\gamma_1, \gamma_2$ respectively, choose $z_1\in \gamma_1$ with $z_1\neq z_0$. By Bell--Kaleem \cite{BeKf}, there exists a proper holomorphic map $f: \Omega \ra \mathbb D$ which maps both $z_1$ and $w_0$ to $1 \in \partial \mathbb D$. We can then shrink $U$, if necessary, to make sure that $f(z)-f(w)$ does not vanish on $\Tilde{U}\times\Tilde{V}$. We also make sure that none of the points $a_i$ lie in $\Tilde{V}$. Then, for $(z,w)\in \Tilde{U}\times \Tilde{V}$, 
		\begin{equation}\label{rel2_k}
			S_{\varphi_k}(z,w) = \frac{1}{1-f(z)\overline{f(w)}} \left(\sum_{i,j=1}^N \sum_{n=0}^{M(i)} \sum_{m=0}^{M(j)} c_{ijnmk}\, S_{\varphi_k}^{(\bar{n})}(z,a_i) \overline{S_{\varphi_k}^{(\bar{m})}(w,a_j)}\right)
		\end{equation}
        where the constants $c_{ijnmk}$ are determined by (\ref{constants_k}). Similarly, for $(z,w)\in \Tilde{U}\times \Tilde{V}$, 
		\begin{equation}\label{rel2}
			S_{\varphi}(z,w) = \frac{1}{1-f(z)\overline{f(w)}} \left(\sum_{i,j=1}^N \sum_{n=0}^{M(i)} \sum_{m=0}^{M(j)} c_{ijnm}\, S_{\varphi}^{(\bar{n})}(z,a_i) \overline{S_{\varphi}^{(\bar{m})}(w,a_j)}\right)
		\end{equation}
        where the constants $c_{ijnm}$ are determined by (\ref{constants}).
        We have seen in the proof of Theorem \ref{Szego} that $c_{ijnmk}\rightarrow c_{ijnm}$ as $k\rightarrow\infty$. We also know that
	\[
	\overline{S_{\varphi_k}(z,w)} = \frac{1}{i\varphi_k(z)}L_{\varphi_k}(z,w)T(z),
	\quad\,\,
	w\in\Om, \,z\in\partial\Omega.
	\]
For a positive integer $r$, upon differentiation with respect to $w$, we obtain
\[
\overline{S_{\varphi_k}^{(\bar{r})}(z,w)} = \frac{1}{i\varphi_k(z)}L_{\varphi_k}^{(r)}(z,w)T(z)
\]
for $w\in\Om$ and $z\in\partial\Omega$. Now (\ref{rel2_k}) can therefore be rewritten for $w\in\partial\Om$, as
\[
\frac{1}{i\varphi_k(w)}L_{\varphi_k}(w,z)T(w) 
=
\frac{1}{1-f(z)\overline{f(w)}} \sum_{i,j=1}^N \sum_{n=0}^{M(i)} \sum_{m=0}^{M(j)} c_{ijnmk}\, S_{\varphi_k}^{(\bar{n})}(z,a_i)
\frac{1}{i\varphi_k(w)} L_{\varphi_k}^{(m)}(w,a_j) T(w).
\]
Since $f(w) \overline{f(w)} = 1$, this can be written as,
\begin{equation}
L_{\varphi_k}(w,z) 
=
\frac{f(w)}{f(w)-f(z)} \sum_{i,j=1}^N \sum_{n=0}^{M(i)} \sum_{m=0}^{M(j)} c_{ijnmk}\, S_{\varphi_k}^{(\bar{n})}(z,a_i)
L_{\varphi_k}^{(m)}(w,a_j), 
\end{equation}
and this holds by the identity principle for all $(z,w)\in \Tilde{U}\times\Tilde{V}$. Similarly,
\begin{equation}
L_{\varphi}(w,z) 
=
\frac{f(w)}{f(w)-f(z)} \sum_{i,j=1}^N \sum_{n=0}^{M(i)} \sum_{m=0}^{M(j)} c_{ijnm}\, S_{\varphi}^{(\bar{n})}(z,a_i)
L_{\varphi}^{(m)}(w,a_j).
\end{equation}
This gives 
\begin{multline*}
l_{\varphi_k}(w,z) - l_{\varphi}(w,z) 
=
L_{\varphi_k}(w,z) - L_{\varphi}(w,z) 
\\= \frac{f(w)}{f(w)-f(z)}
 \sum_{i,j=1}^N \sum_{n=0}^{M(i)} \sum_{m=0}^{M(j)} 
 c_{ijnmk}\, S_{\varphi_k}^{(\bar{n})}(z,a_i)
L_{\varphi_k}^{(0,m)}(w,a_j)
-
c_{ijnm}\, S_{\varphi}^{(\bar{n})}(z,a_i)
L_{\varphi}^{(0,m)}(w,a_j).
\end{multline*}
for $(z,w)\in \Tilde{U}\times\Tilde{V}$
Note that $L_{\varphi_k}(w,z) - L_{\varphi}(w,z)$ is a holomorphic function in $z$ and $w$ and extends to be in $C^{\infty}(\overline{\Omega}\times \overline{\Omega})$. The function on the right hand side is also a holomorphic function for $(z,w)\in\Tilde{U}\times\Tilde{V}$ as none of the points $a_i$ lie in $\Tilde{V}$. 
By \cite{JV}, 
\[
    \lim_{k\rightarrow\infty}S_{\varphi_k}^{(\bar{n})}(z,a_i)= S_{\varphi}^{(\bar{n})}(z,a_i)
	\quad
	\text{and}
	\quad
	\lim_{k\rightarrow\infty}L_{\varphi_k}^{(0,m)}(w,a_j)= L_{\varphi}^{(0,m)}(w,a_j)
\]
uniformly for $z \in \Tilde{U}$ and $w \in \Tilde{V}$. Since the constants $c_{ijnmk}$ converge to $c_{ijnm}$ as $k\rightarrow\infty$, it follows that 
\[
		\lim_{k\rightarrow\infty}l_{\varphi_k}(z,w) = l_{\varphi}(z,w)
		\]
		locally uniformly on $(\overline{\Om}\times\overline{\Om})\setminus\Delta$. Note that the convergence of derivatives also follows similarly. 
\end{proof}
	
	
	%%%%%%%%%%%%%%%%%%%%%%%%%%%%%%%%%%%%
	
		
	\section{zeros of the weighted Szeg\H{o} kernel}
	
\noindent In the unweighted case, the following behaviour of the zeros of the Szeg\H{o} kernel is known and an exposition can be found in Bell \cite{Be0}.
	
	\begin{thm}[Bell]
		Let $\Omega\subset\mathbb{C}$ be a bounded $n$-connected domain with $C^{\infty}$ smooth boundary curves $\gamma_i$, $1 \le i \le n$. Let $w_j$ be a sequence in $\Omega$ that tends to a point $a \in \gamma_m$. As $w_j$ tends to $a$, the zeros $Z^i(w_j)$ of $S(z,w_j)$ become simple zeros, and it is possible to order them so that for each $i$, $i\neq m$, there is a point $Z^i(a)\in\gamma_i$ such that $Z^i(w_j)$ tends to $Z^i(a)$. 
		
		\medskip
		
		$S(z,a)$ is non-vanishing on $\Om$ as a function of $z$ and has exactly $n-1$ zeros $Z^i(a)$ on $\partial\Om$, one on each boundary component not containing $a$. Furthermore, the zeros are simple in the sense that $S'(Z^i(a),a)\neq 0$ for each zero $Z^i(a)$. 
	\end{thm}
	
	\begin{rem}
		It is known that $L(z, w)$ does not vanish for $z\in \overline{\Om}, w\in\Om$. From above theorem, given $a\in\partial\Om$, $L(z,a)$ is non-vanishing for $z\in\Om$ and has exactly $n-1$ zeros on $\partial\Om$, which are same as the zeros of $S(z,a)$.
	\end{rem}
	
	We shall now study the behaviour of zeros of the weighted Szeg\H{o} kernel and the weighted Garabedian kernel as weights vary.
	
	\medskip
	
	Let $\Omega\subset\mathbb{C}$ be a bounded $n$-connected domain whose boundary curves are {\it real analytic} with no isolated boundary point. If $\varphi$ is a positive real analytic function on $\partial\Omega$, then the weighted Szeg\H{o} kernel $S_{\varphi}(z,w)$ extends to a neighborhood of $(\overline{\Om}\times \overline{\Om})\setminus\Delta$ as a function that is holomorphic in $z$ and antiholomorphic in $w$, and the weighted Garabedian kernel is given by
	\[
	L_{\varphi}(z,w) = \frac{1}{2\pi(z-w)}+ l_{\varphi}(z,w)
	\]
	where $l(z,w)$ extends to be holomorphic in $z$ and $w$ on a neighborhood of ($\overline{\Om}\times\ov{\Om})$. Fix $w\in\Om$. Let $z_1,\ldots,z_N$ be the zeros of $S_{\varphi}(\cdot,w)$ in $\Om$ and $b_1,\ldots,b_M$ be the zeros of $S_{\varphi}(\cdot, w)$ on $\partial\Om$. We know that
	\begin{equation}\label{Bdy}
		\overline{S_{\varphi}(z,w)} = \frac{1}{i\varphi(z)}L_{\varphi}(z,w)T(z),
		\quad\,\,
		w\in\Om, z\in\partial\Omega.
	\end{equation}
	Therefore, the zeros of $L_{\varphi}(\cdot, w)$ on $\partial\Om$ are $b_1,\ldots, b_M$. Let $p_1,\ldots, p_Q$ be the zeros of $L_{\varphi}(\cdot,w)$ in $\Om$. We shall apply the generalized argument principle on $h = S_{\varphi}(\cdot,w)L_{\varphi}(\cdot,w)$, which is a meromorphic function in a neighborhood of $\overline{\Om}$. Let $m_h(z)$ denote the multiplicity of a zero or pole of $h$ at $z$. We obtain
	\[
	\sum_{i=1}^N m_h(z_i) + \sum_{i=1}^Q m_h(p_i) - m_h(w) + \frac{1}{2}\sum_{i=1}^M m_h(b_i)
	=
	\frac{1}{2\pi} \Delta \arg h.
	\]
	Let $m_S(z)$ and $m_L(z)$ denote the the multiplicity of a zero or a pole of $S_{\varphi}(\cdot,w)$ and $L_{\varphi}(\cdot, w)$ respectively at $z$. Then,
	\[
	\sum_{i=1}^N m_S(z_i) + \sum_{i=1}^Q m_L(p_i) - m_L(w) + \sum_{i=1}^M m_S(b_i) 
	=
	\frac{1}{2\pi} \Delta \arg S_{\varphi}(\cdot,w)L_{\varphi}(\cdot,w).
	\]
	Multiplying $S_{\varphi}(z,w)$ on both sides of (\ref{Bdy}), we obtain
	\[
	\frac{1}{i\varphi(z)}S_{\varphi}(z, w) L_{\varphi}(z,w) T(z)= \vert S_{\varphi}(z,w)\vert^2.
	\]
	Therefore, 
	\[
	\Delta \arg S_{\varphi}(\cdot,w)L_{\varphi}(\cdot,w) = - \Delta \arg T = 2\pi(n-2).
	\]
	Since $m_L(w) = 1$, we have
	\begin{equation}
	\sum_{i=1}^N m_S(z_i) + \sum_{i=1}^Q m_L(p_i) + \sum_{i=1}^M m_S(b_i) 
	=
	n-1.
	\end{equation}
	\begin{obs}
		Thus, for $w\in\Omega$, the combined number of zeros of $S_{\varphi}(\cdot, w)$ and $L_{\varphi}(\cdot, w)$ in $\Om$, plus the number of zeros of $S_{\varphi}(\cdot,w)$ on $\partial\Om$ (which are also the zeros of $L_{\varphi}(\cdot,w)$ on $\partial\Om$), equals $n-1$. Furthermore, the zeros are counted with multiplicity.
	\end{obs}
	
	\medskip
	
	Let $\{\varphi_k\}_{k=1}^{\infty}$ be a sequence of positive real analytic functions on $\partial\Omega$ such that $\varphi_k\rightarrow 1$ in $C^{\infty}$ topology on $\partial\Omega$ as $k\rightarrow\infty$. Let $a\in\partial\Omega$. 
	
	\begin{thm}	
		There exists $k_0 \ge 1$ such that for all $k\geq k_0$ the following holds: 
		
		\medskip
		
		Let $w_j\in\Om$ be a sequence that converges to $a$. The functions $S_{\varphi_k}(\cdot,w_j)$ and $L_{\varphi_k}(\cdot,w_j)$ have distinct zeros in $\Omega$ for large enough $j$. The combined zeros $Z^i_k(w_j)$ of $S_{\varphi_k}(z,w_j)$ and $L_{\varphi_k}(\cdot, w_j)$ in $\overline{\Omega}$, become simple zeros, and it is possible to order them so that for each $i=1,\ldots, n-1$, there is a point $Z_k^i(a)\in\ov{\Om}$ such that $Z^i_k(w_j)$ tends to $Z^i_k(a)$ as $j\rightarrow\infty$. 
		
		\medskip
		
		The functions $S_{\varphi_k}(\cdot,a)$ and $L_{\varphi_k}(\cdot,a)$ have distinct zeros in $\Omega$. The combined number of zeros of $S_{\varphi_k}(\cdot, a)$ and $L_{\varphi_k}(\cdot,a)$ in $\Om$, plus the number of zeros of  $S_{\varphi_k}(\cdot,a)$ on $\partial\Om$ $($which are also the zeros of $L_{\varphi_k}(\cdot,a)$ on $\partial\Om)$, counting multiplicity, equals $n-1$, and are same as $Z_{k}^i(a)$. Furthermore, all $Z_{k}^i(a)$ are distinct for $i=1,\ldots, n-1$, and therefore the zeros are simple.
		
		\medskip
		
		Moreover, 
		\[
		\lim_{j,k\rightarrow\infty} Z^i_{k}(w_j) = Z^i(a)
		\]
		where $Z^i(a)$ are the zeros of $S(\cdot,a)$ (and $L(\cdot,a)$ as well).
	\end{thm}
	
	\begin{proof}
		Let $\gamma_i$ denote the $n$ boundary curves of $\Om$. For the sake of simplicity, we denote the boundary curves in such a manner that $a\in \gamma_n$. 
		Let $\Hat{\Om}$ denote the double of $\Om$, $R(z)$ denote the antiholomorphic involution on $\Hat{\Om}$ which fixes the boundary, and $\Tilde{\Om} = R(\Om)$ denote the reflection of $\Om$ in $\Hat{\Om}$ across the boundary. For each $i\neq n$, choose a smoothly bounded simply connected neighborhood $U_i$ of $Z^i(a)$ in $\Hat{\Om}$ and a point $b\in \Om$ such that 
		
		\begin{enumerate}
			\item[(i)] $\overline{U_i}$ are mutually disjoint
			\item[(ii)] $z\in U_i$ if and only if $R(z)\in U_i$
			\item[(iii)] $\ov{U_i}\cap ({\gamma_n \cup \{b\}}) = \emptyset$ 
			\item[(iv)] $S(z,b) \neq 0$ for all $z\in \overline{U_i}\cap \ov{\Om}$.
		\end{enumerate}
		
		\medskip
		
		Let $C_i$ denote the boundary curve of $U_i$. We claim that there exists a neighborhood $U$ of $a$ in $\overline{\Om}$ such that $S(z,w)\neq 0$ for all $z\in C_i\cap \overline{\Om}$ and $w\in \ov{U}$. If it were not true, then for every positive integer $m$, there would exist $w_m\in\ov{\Om}$ with $\vert w_m - a\vert \leq 1/m$ and $z_m\in C_i\cap \overline{\Om}$ for some $i$ such that $S(z_m,w_m) = 0$. By the pigeonhole principle, infinitely many $z_m$ lie on some $C_{i_0}$. So, there exists a subsequence $z_{m_r}$ with $z_{m_r}\rightarrow z_0\in C_{i_0}$ as $r\rightarrow\infty$. This implies that $S(z_0,a) = 0$ which is a contradiction.
		
		\medskip
		
		We know that
		\begin{equation}\label{thms}
			\lim_{k\rightarrow\infty} S_{\varphi_k}(z,w) = S(z,w)
			\quad \text{and}\quad
			\lim_{k\rightarrow\infty} L_{\varphi_k}(z,w) = L(z,w)
		\end{equation}
		locally uniformly on $(\ov{\Om}\times\ov{\Om})\setminus \{(z,z): z\in\partial\Om\}$ and $(\ov{\Om}\times\ov{\Om})\setminus \{(z,z): z\in \overline{\Om}\}\}$ respectively. Therefore, we can choose an integer $k_0 \ge 1$ such that for all $k\geq k_0$, 
		
		\begin{enumerate}
			\item[(i)] $S_{\varphi_k}(z,w)$ does not vanish for $z\in C_i\cap \overline{\Om}$, $w\in \ov{U}$
			\item[(ii)] $L_{\varphi_k}(R(z),w)$ does not vanish for $z\in C_i\cap (\Tilde{\Om}\cup \partial\Om)$, $w\in \ov{U}$
			\item[(iii)] $S_{\varphi_k}(z,b)$ does not vanish for $z\in \ov{U_i}\cap \overline{\Om}$
			\item[(iv)] $L_{\varphi_k}(R(z),b)$ does not vanish for $z\in \ov{U_i}\cap (\Tilde{\Om}\cup \partial\Om)$.
		\end{enumerate}
		
		\medskip
		
		For $k\geq k_0$ and $w\in \ov{U}$, define the functions $G_k(z,w)$ by
		\begin{equation}
		G_k(z,w)
		=
		\begin{cases}
			\mathcal{S}_{\varphi_k}(z,w)& z\in U_i\cap \overline{\Omega}
			\\
			\overline{\mathcal{L}_{\varphi_k}(R(z),w)}& z\in U_i\cap \Tilde{\Omega}
		\end{cases}
		\end{equation}
		where 
		\begin{equation}
		\mathcal{S}_{\varphi_k}(z,w) = \frac{S_{\varphi_k}(z,w)}{S_{\varphi_k}(z,b)}
		\quad
		\text{and}
		\quad
		\mathcal{L}_{\varphi_k}(z,w) = \frac{L_{\varphi_k}(z,w)}{L_{\varphi_k}(z,b)}.
		\end{equation}
		Recall that
		\[
			\overline{S_{\varphi_k}(z,w)} = \frac{1}{i\varphi_k(z)}L_{\varphi_k}(z,w)T(z),
			\quad\,\,
			w\in\Om, z\in\partial\Omega.
		\]
		This implies that
		\[
		\mathcal{S}_{\varphi_k}(z,w)
		=
		\frac{S_{\varphi_k}(z,w)}{S_{\varphi_k}(z,b)}
		=
		\overline{\left(\frac{L_{\varphi_k}(z,w)}{L_{\varphi_k}(z,b)}\right)}
		=
		\overline{\left(\frac{L_{\varphi_k}(R(z),w)}{L_{\varphi_k}(R(z),b)}\right)}
		=
		\overline{\mathcal{L}_{\varphi_k}(R(z),w)},
		\quad
		z\in\partial\Om.
		\]
		Therefore, the functions $G_k$ are well-defined and holomorphic on $U_i$ and extend smoothly to $C_i$. Similarly, we define the function $G(z,w)$ by
		\begin{equation}
		G(z,w)
		=
		\begin{cases}
			\mathcal{S}(z,w)& z\in U_i\cap \overline{\Omega}
			\\
			\overline{\mathcal{L}(R(z),w)}& z\in U_i\cap \Tilde{\Omega}
		\end{cases}
		\end{equation}
		where 
		\begin{equation}
		\mathcal{S}(z,w) = \frac{S(z,w)}{S(z,b)}
		\quad
		\text{and}
		\quad
		\mathcal{L}(z,w) = \frac{L(z,w)}{L(z,b)}.
		\end{equation}
		The function $G$ is well-defined, holomorphic on $U_i$ and extends smoothly to $C_i$. 
		Now, the function 
		\begin{equation}
		N_{ki}(w) = \int_{C_i} \frac{{G_k}'(z, w)}{G_k(z,w)} dz
		\end{equation}
		converges uniformly to 
		\begin{equation}
		N_i(w)=\int_{C_i}  \frac{{G}'(z, w)}{G(z,w)} dz
		\end{equation}
		for $w\in \overline{U}$ as $k\rightarrow\infty$. Here, $N_{ki}(w)$ and $N_i(w)$ are the number of zeros of $G_k(\cdot, w)$ and $G(\cdot,w)$ in $U_i$, respectively. We go back and choose a larger $k_0 \ge 1$, if needed, so that
		\[
		\vert N_{ki}(w) - N_i(w)\vert < 1/2
        \]
		for every $i,w\in \overline{U}$ and $k\geq k_0$.

        \medskip
		
		We shall now fix a $k\geq k_0$. Let $w_j\in\Omega$ be a sequence that converges to $a$. Assume, without loss of generality, that $w_j\in U$ for all $j$. The zeros of $G(\cdot, w_j)$ in $U_i$ are same as the zeros of $S(\cdot, w_j)$. So, there exists $j_0 \ge 1$ such that $N_i(w_j) = 1$ for all $i$ and $j\geq j_0$. Therefore, $N_{ki}(w_j) = 1$ for all $i$ and $j\geq j_0$. By the continuity of $N_{ki}(w)$ as a function of $w$, this implies that $N_{ki}(a) = 1$.
		
		\medskip
		
		The zeros of $G_k(\cdot,w_j)$ are the zeros of $S_{\varphi_k}(\cdot,w_j)$ or $L_{\varphi_k}(R(\cdot),w_j)$. The combined number of zeros of $S_{\varphi_k}(\cdot, w_j)$ and $L_{\varphi_k}(\cdot, w_j)$ in $\Om$, plus the number of zeros of $S_{\varphi_k}(\cdot,w_j)$ on $\partial\Om$ (which are also the zeros of $L_{\varphi_k}(\cdot,w_j)$ on $\partial\Om$), equals $n-1$. Since $N_{ki}(w_j) =1$ for all $i=1,\ldots,n-1$ and $j\geq j_0$, the functions $S_{\varphi_k}(\cdot,w_j)$ and $L_{\varphi_k}(\cdot,w_j)$ must have distinct zeros in $\Om$. Furthermore, the combined zeros $Z^i_k(w_j)$ of $S_{\varphi_k}(\cdot,w_j)$ and $L_{\varphi_k}(\cdot,w_j)$ in $\overline{\Om}$, are simple for all $j\geq j_0$. We order these zeros in such a manner that $Z^i_k(w_j)\in U_i$ for all $j\geq j_0$.
		
		\medskip
		
		Let $Z_i(G_k(\cdot,w_j))$ and $Z_i(G_k(\cdot,a))$ denote the zero of $G_k(\cdot,w_j)$ and $G_k(\cdot,a)$ in $U_i$ respectively. By an application of the residue theorem and (\ref{thms}),
		\[
		Z_i(G_k(\cdot,w_j))= \int_{C_i} \frac{z\,{G_k}'(z, w_j)}{G_k(z,w_j)} dz 
		\quad 
		\rightarrow
		\quad
		\int_{C_i}  \frac{z\,{G_k}'(z, a)}{G_k(z,a)} dz 
		=
		Z_i(G_k(\cdot,a))
        \]
		as $j\rightarrow\infty$.

    \medskip
        
		If $Z_i(G_k(\cdot,a))\in \ov{\Om}$, we let $Z^i_k(a) = Z_i(G_k(\cdot,a))$. Otherwise, we let $Z^i_k(a) = R(Z_i(G_k(\cdot,a)))$. Then, 
		\[
		\lvert Z^{i}_k(w_j)- Z^i_k(a) \vert 
		\leq 
		\vert Z_i(G_k(\cdot,w_j))- Z_i(G_k(\cdot,a)) \vert
		\rightarrow 0
		\quad
		\text{as }j\rightarrow\infty.
		\]
		Furthermore, $Z^i_k(a)\in U_i$ are distinct.
		
		\medskip
		
		Since the zeros of $G_k(\cdot,a)$ are the zeros of $S_{\varphi_k}(\cdot,a)$ or $L_{\varphi_k}(R(\cdot),a)$, the points $Z^i_k(a)$ are the zeros of $S_{\varphi_k}(\cdot,a)$ or $L_{\varphi_k}(\cdot,a)$ in $\overline{\Om}$. Let, if possible, $z_0\in\ov{\Om}$ be another zero of $S_{\varphi_k}(\cdot,a)$ or $L_{\varphi_k}(\cdot,a)$. Let us assume, without loss of generality, that $S_{\varphi_k}(z_0,a) = 0$, as a similar reasoning can be applied if $L_{\varphi_k}(z_0,a)=0$.  It is clear that $z_0\notin \overline{(U_i\cap \Om)}$ for any $i=1,\ldots,n-1$. If $z_0\in\Om$, choose a ball $B(z_0,\delta)\subset\subset\Om$ centered at $z_0$ of radius $\delta>0$ such that 
		
		\begin{itemize}
			\item[(i)] $\overline{B(z_0,\delta)}$ is disjoint from $\ov{U_i\cap \Om}$
			\item[(ii)] $S_{\varphi_k}(\cdot,a)$ is non-vanishing on $C=\partial B(z_0.\delta)$. 
		\end{itemize}
		
		\noindent Then $S_{\varphi_k}(\cdot,w_j)$ is also non-vanishing on $C$ for all large $j$. If it were not true, then for every $j \ge 1$, there would exist $z_j\in C$ such that $S_{\varphi_k}(z_j,w_j) = 0$. By passing to a convergent subsequence of $z_j$, we would obtain a zero of $S_{\varphi_k}(\cdot,a)$ on $C$, which is not possible. As $j\rightarrow\infty$,
		\[
		\int_{C}\frac{{S_{\varphi_k}}'(z,w_j)}{S_{\varphi_k}(z,w_j)} dz
		\quad
		\rightarrow
		\quad
		\int_{C}\frac{{S_{\varphi_k}}'(z,a)}{S_{\varphi_k}(z,a)} dz
		\quad
		\geq 1.
		\]
		So, there must exist a zero of $S_{\varphi_k}(\cdot,w_j)$ in $B(z_0,\delta)$ for all large $j$, hence a contradiction. If $z_0\in\partial\Om$, we will work with $G_k(\cdot,a)$ on an appropriate neighborhood of $z_0$, and arrive at a contradiction by a similar reasoning. Thus, we have proved that $Z^i_k(a)$ are the only zeros of $S_{\varphi_k}(\cdot,a)$ and $L_{\varphi_k}(\cdot,a)$. Since $N_{ki}(a) = 1$, it follows that $S_{\varphi_k}(\cdot,a)$ and $L_{\varphi_k}(\cdot,a)$ have distinct zeros in $\Om$. 
		
		\medskip
		
		We have also proved that the combined number of zeros of $S_{\varphi_k}(\cdot, a)$ and $L_{\varphi_k}(\cdot,a)$ in $\Om$, plus the number of zeros of  $S_{\varphi_k}(\cdot,a)$ on $\partial\Om$ $($which are also the zeros of $L_{\varphi_k}(\cdot,a)$ on $\partial\Om)$, counting multiplicity, equals $n-1$, and are same as $Z_{k}^i(a)$. Moreover, these zeros are simple.
		
		\medskip
		
		Finally, let $\epsilon>0$ be given. Using (\ref{thms}), we can choose $k_1\ge 1$ with $k_1\geq k_0$ such that
		\[
		\left\lvert
		\int_{C_i} z \; \frac{{G'_k}(z, w)}{G_k(z,w)} dz - \int_{C_i}  z \; \frac{G'(z, w)}{G(z,w)} dz 
		\right\rvert
		\,\,<\,\,
		\frac{\epsilon}{2}
		\quad
		\text{for all }i=1,\ldots, n-1,
		\,w\in\overline{U}
		\text{ and }
		k\geq k_1.
		\]
		Choose $j_1 \ge 1$ with $j_1\geq j_0$ such that
		\[
		\left\lvert
		\int_{C_i} z \; \frac{G'(z, w_j)}{G(z,w_j)} dz - \int_{C_i} z \; \frac{G'(z, a)}{G(z,a)} dz 
		\right\rvert 
		\,\,<\,\,
		\frac{\epsilon}{2}
		\quad
		\text{for all }i=1,\ldots, n-1
		\text{ and }
		j\geq j_1.
		\]
		Since $Z^i(a)\in\partial\Omega$, we have $\vert Z^i_k(w_j) - Z^i(a) \vert = \vert Z_i(G_k(\cdot,w_j)) - Z^i(a)\vert$. Therefore, for $j\geq j_1$ and $k\geq k_1$, we obtain using the residue theorem that
		\begin{multline*}
			\vert Z^i_k(w_j) - Z^i(a) \vert
			=
			\vert Z_i(G_k(\cdot,w_j)) - Z^i(a)\vert
			\\
			=
			\left\lvert
			\int_{C_i} z \; \frac{G'_k(z, w_j)}{G_k(z,w_j)} dz - \int_{C_i} z \; \frac{G'(z, a)}{G(z,a)} dz
			\right\rvert
			\\
			\leq
			\left\lvert
			\int_{C_i} z \; \frac{G'_k(z, w_j)}{G_k(z,w_j)} dz - \int_{C_i} z \;\frac{G'(z, w_j)}{G(z,w_j)} dz
			\right\rvert
			+
			\left\lvert
			\int_{C_i} z \;\frac{G'(z, w_j)}{G(z,w_j)} dz - \int_{C_i} z \; \frac{G'(z, a)}{G(z,a)} dz
			\right\rvert
			\\
			<
			\frac{\epsilon}{2} + \frac{\epsilon}{2} 
			= \epsilon.
		\end{multline*}
		Thus, 
		\begin{equation}
		\lim_{j,k\rightarrow\infty} Z^i_{k}(w_j) = Z^i(a)
		\end{equation}
		and this completes the proof of the theorem.		
	\end{proof}
	
	
	%%%%%%%%%%%%%%%%%%%%%%%%%%%%%%%%%%%5
	
	
	\section{Weighted Ahlfors maps}
	
	\begin{thm}[Nehari \cite{Nehari}]
		Let $\Omega\subset\mathbb{C}$ be a bounded $n$-connected domain with $C^{\infty}$ smooth boundary such that no boundary curve is an isolated point. Let $\varphi$ be positive real-valued $C^{\infty}$ smooth function on $\partial\Om$. Given $a\in\Om$, let
		\[
		\mathcal{B}_{\varphi} = \left\lbrace
		f\in\mathcal{O}(\Om) : f(a) = 0 \text{ and } \limsup_{z\rightarrow z_0}\vert f(z)\vert \leq \frac{1}{\varphi(z_0)}\text{ for }z_0\in\partial\Om\right\rbrace.
		\]
		Then for every $f\in \mathcal{B}_{\varphi}$,
		\begin{equation}\label{ahlfors}
			\vert f'(a)\vert \leq \vert F'(a) \vert = 2\pi S_{\varphi}(a,a)
		\end{equation}
		where 
		\begin{equation}
		F(z) = \frac{S_{\varphi}(z,a)}{L_{\varphi}(z,a)}
		\end{equation}
		is the weighted Ahlfors map. If $L_{\varphi}(z,a) \neq 0$ for all $z\in \Om$, then $F\in \mathcal{B}_{\varphi}$ and it is the unique function (modulo rotation) in $\mathcal{B}_{\varphi}$ satisfying (\ref{ahlfors}).
	\end{thm}
	

		Since 
		\[
		\overline{S_{\varphi}(z,a)} = \frac{1}{i\varphi(z)} L_{\varphi}(z,a) T(z), \quad z\in\partial\Om,
		\]
		the $\lim_{z\rightarrow z_0}\vert F(z)\vert$ exists for all $z_0\in\partial\Om$, and
		\begin{equation}
		\lim_{z\rightarrow z_0} \vert F(z)\vert 
		=
		\frac{1}{\varphi(z_0)},\quad z_0\in\partial\Om.
		\end{equation}
	Thus, $F$ has a well defined limit on $\pa \Omega$ and this strengthens Nehari's theorem mentioned above.
	
	\begin{thm}
		Let $\Omega\subset\mathbb{C}$ be a bounded $n$-connected domain with $C^{\infty}$ smooth boundary such that no boundary curve is an isolated point. Let $\{\varphi_k\}_{k=1}^{\infty}$ be a sequence of positive real-valued $C^{\infty}$ smooth functions on $\partial\Omega$ such that $\varphi_k\rightarrow 1$ in the $C^{\infty}$ topology on $\partial\Omega$ as $k\rightarrow\infty$. For a compact set $W\subset\Omega$, there exists $k_0 \ge 1$ such that for all $a\in W$ and $k\geq k_0$
		\begin{itemize}
			\item[(i)] $L_{\varphi_k}(\cdot,a)$ does not vanish on $\overline{\Om}$ 
			\item[(ii)] $S_{\varphi_k}(\cdot,a)$ has $n-1$ zeros in $\Om$ and does not vanish on $\partial\Om$ 
		\end{itemize}
		Furthermore, the weighted Ahlfors maps with respect to $a\in W$ converge to the corresponding classical Ahlfors map, uniformly on $\ov{\Om}$. The convergence is also uniform with respect to $a\in W$. In other words,
		\begin{equation}
		\lim_{k\rightarrow\infty} F_k(z,a)
		:=
		\lim_{k\rightarrow\infty} \frac{S_{\varphi_k}(z,a)}{L_{\varphi_k}(z,a)} 
		=
		\frac{S(z,a)}{L(z,a)} 
		=:
		F(z,a)
		\end{equation}
		uniformly for $a\in W$ and $z\in\overline{\Om}$.
	\end{thm}
	
	\begin{proof}
		Choose $k_1 \ge 1$ such that $L_{\varphi_k}(z,a)$ and $S_{\varphi_k}(z,a)$ do not vanish for $k\geq k_1$, $z\in \partial\Om$ and $a\in W$. For $k\geq k_1$, 
		\[
		\frac{1}{i\varphi_k(z)}S_{\varphi_k}(z,a)L_{\varphi_k}(z,a) T(z) = \vert S_{\varphi_k}(z,a)\vert^2,
		\quad z\in\partial\Om, a\in\Om
		\]
		and this gives
		\[
		\Delta \arg S_{\varphi_k}(\cdot, a)L_{\varphi_k}(\cdot, a) = -\Delta \arg T = 2\pi (n-2)
		\]
		Since $L_{\varphi_k}(z,a)$ has a simple pole at $z=a$, the argument principle implies that the combined number of zeros of $S_{\varphi_k}(\cdot,a)$ and $L_{\varphi_k}(\cdot,a)$ in $\Om$ equals $n-1$, counting multiplicity, for all $a\in W$. The function
		\[
		N_k(a) = \int_{\partial \Om}\frac{{S_{\varphi_k}}'(z,a)}{S_{\varphi_k}(z,a)} dz
		\]
		gives total number of zeros of $S_{\varphi_k}(\cdot,a)$ in $\Om$. The function $N_k$ converges uniformly for $a\in W$ to the function
		\[
		N(a) =  \int_{\partial \Om}\frac{S'(z,a)}{S(z,a)} dz
		\]
		that gives the number of zeros of $S(\cdot,a)$ in $\Om$, counting multiplicities, which is $n-1$. Therefore, eventually $N_k$ are constant functions. Choose $k_0\geq k_1$ such that $N_k(a) = n-1$ for all $k\geq k_0$ and $a\in W$. Thus, $S_{\varphi_k}(\cdot,a)$ has $n-1$ zeros in $\Om$ and does not vanish on $\partial\Om$ for all $k\geq k_0$ and $a\in W$. Since the combined number of zeros of $S_{\varphi_k}(\cdot,a)$ and $L_{\varphi_k}(\cdot,a)$ equals $n-1$, the function $L_{\varphi_k}(\cdot,a)$ does not vanish on $\overline{\Om}$ for all $k\geq k_0$ and $a\in W$.
		
		\medskip
		
		We know that
		\[
		\lim_{k\rightarrow\infty} S_{\varphi_k}(z,w) = S(z,w)
		\quad \text{and}\quad
		\lim_{k\rightarrow\infty} L_{\varphi_k}(z,w) = L(z,w)
		\]
		locally uniformly on $(\ov{\Om}\times\ov{\Om})\setminus \{(z,z): z\in\partial\Om\}$ and $(\ov{\Om}\times\ov{\Om})\setminus \{(z,z): z\in \overline{\Om}\}\}$ respectively. Therefore,
		\[
		\lim_{k\rightarrow\infty} F_k(z,a) = F(z,a)
		\]
		uniformly for $a\in W$ and $z\in\partial\Om$. Since $F_k(\cdot,a),F(\cdot,a)\in A^{\infty}(\Om)$ for all $k\geq k_0$ and $a\in W$, it follows by the maximum modulus principle that 
		\begin{eqnarray*}
			\sup\{ \vert F_k(z,a) - F(z,a)\vert : z \in \overline{\Om}, a\in W\} 
			&=&
			\sup\{ \vert F_k(z,a) - F(z,a)\vert : z \in \partial\Om, a\in W\}.
		\end{eqnarray*}
		Thus proves that $\lim_{k\rightarrow\infty} F_k(z,a) = F(z,a)$ uniformly for $a\in W$ and $z\in\overline{\Om}$.
	\end{proof}


    %%%%%%%%%%%%%%%%%%%%%%%%%%
    \section{Explicit Formulas for a class of weights}
	
	We start with two examples due to Nehari and Bell, giving more details.
	
	\begin{exam}[Nehari \cite{Nehari}]
		Let $\Omega\subset\mathbb{C}$ be a bounded $n$-connected domain with $C^{\infty}$ smooth boundary such that no boundary curve is an isolated point. For $a\in\Om$, let $a_1,\ldots a_r$ be the zeros of $S(z,a)$ with multiplicity $n_1,\ldots n_r$, respectively. For $1\leq k\leq r$ and $1\leq m_i\leq n_i$, let
		\[
		\varphi(z) = \vert (z - a_1)^{m_1} \ldots (z-a_k)^{m_k} \vert^2.
		\]
		Given $g\in H^2(\partial\Omega)$,
		\begin{eqnarray*}
			g(a) 
			&=&
			\frac{1}{\prod_{i=1}^k (a-a_i)^{m_i}} \, \left( g(a) \prod_{i=1}^k (a-a_i)^{m_i} \right)
			\\&=&
			\frac{1}{\prod_{i=1}^k (a-a_i)^{m_i}} 
			\int_{\partial\Om} g(\zeta) \prod_{i=1}^k (\zeta-a_i)^{m_i} \,\overline{S(\zeta,a)} \,ds
			\\&=&
			\frac{1}{\prod_{i=1}^k (a-a_i)^{m_i}} 
			\int_{\partial\Om} g(\zeta) \prod_{i=1}^k (\zeta-a_i)^{m_i} \,\overline{S(\zeta,a)} \frac{\varphi(\zeta)}{\prod_{i=1}^k \vert \zeta - a_i\vert^{2 m_i}}\,ds
			\\&=&
			\int_{\partial\Om} g(\zeta) \,\frac{\overline{S(\zeta,a)}}{\prod_{i=1}^k (\overline{\zeta} - \overline{a_i})^{m_i}(a-a_i)^{m_i}} \,\varphi(\zeta)\,ds.
		\end{eqnarray*}
		Therefore, by the uniqueness of the weighted Szeg\H{o} kernel function $S_{\varphi}(\cdot, a)$, we have
		\[
		S_{\varphi}(z,a) = \frac{S(z,a)}{\prod_{i=1}^k (z-a_i)^{m_i} (\overline{a}-\overline{a_i})^{m_i}}.
		\]
		Now for $z\in\partial\Omega$,
		\begin{eqnarray*}
			L_{\varphi}(z, a) 
			&=&
			i \varphi(z) \overline{T(z)} \overline{S_{\varphi}(z,a)} 
			\\&=&
			i \prod_{i=1}^{k} \vert z-a_i\vert^{2m_i} \,\overline{T(z)}  \frac{\overline{S(z,a)}}{\prod_{i=1}^k (\overline{z}-\overline{a_i})^{m_i} (a-a_i)^{m_i}}
			\\&=&
			\prod_{i=1}^{k}\frac{(z-a_i)^{m_i}}{(a-a_i)^{m_i}} \,L(z,a).
		\end{eqnarray*}
		The identity principle implies that
		\[
		L_{\varphi}(z, a) = \prod_{i=1}^{k}\frac{(z-a_i)^{m_i}}{(a-a_i)^{m_i}} \,L(z,a), \quad z\in\overline{\Om}.
		\]
		So, the weighted Garabedian kernel function $L_{\varphi}(\cdot,a)$ does not vanish on $\partial\Om$ and has zeros $a_1,\ldots, a_k$ in $\Om$ with multiplicity $m_1,\ldots, m_k$ respectively. 
		
		\medskip
		
		The weighted Ahlfors map $F_{\varphi}$ with respect to $a\in\Om$, is given by
		\[
		F_{\varphi}(z) = \frac{S_{\varphi}(z,a)}{L_{\varphi}(z,a)} 
		=
		\prod_{i=1}^k\frac{(a-a_i)^{2m_i}}{\vert a - a_i\vert^{2m_i}} \frac{1}{(z-a_i)^{2m_i}} 
		\frac{S(z,a)}{L(z,a)}.
		\]
		Thus, the weighted Ahlfors map $F_{\varphi}$ and the classical Ahlfors map $F$ with respect to $a\in \Om$ are related by
		\[
		F_{\varphi}(z) = \prod_{i=1}^k\frac{(a-a_i)^{2m_i}}{\vert a - a_i\vert^{2m_i}} \frac{1}{(z-a_i)^{2m_i}} 
		F(z).
		\] 
		If $2m_i > n_i$, then the function $F_{\varphi}$ has a pole at $a_i$ of order $2m_i-n_i$. If $2m_i\leq n_i$ for each $i$, then $F_{\varphi}\in \mathcal{B}_{\varphi}$.
		
		\medskip
		
		The function $F_{\varphi}$ maps the $n$ boundary curves $\gamma_i$ of $\Om$ onto $n$ Jordan closed curves $\Gamma_i$. Some of the images $\Gamma_i$ may coincide, for example all of them coincide if $\varphi\equiv 1$. The image of $\Om$ under $F_{\varphi}$ is open and connected and the boundary is the union of $\Gamma_i$. Since $F_{\varphi}$ has a pole, $F_{\varphi}(\Om)$ is the connected component of 
		\[
		(\mathbb{C}\cup\{\infty\} )\setminus \cup_{i=1}^n\Gamma_i
		\]
		that contains $\infty$.
	\end{exam}
	
	\begin{exam}[Bell \cite{Be1}]
		Let $\Omega\subset\mathbb{C}$ be a bounded $n$-connected domain with $C^{\infty}$ smooth boundary such that no boundary curve is an isolated point. Let $p(a,z)$ denote the Poisson kernel of $\Om$. For $A_0\in\Om$, let 
		\[
		\varphi(z) = p(A_0,z).
		\]
		Note that the weight $\varphi$ is well-defined for simply connected domains as well. For $g\in H^2(\partial\Om)$,
		\[
		g(A_0) = \int_{\partial\Om} g(\zeta) \varphi(\zeta) ds,
		\]
		by the reproducing property of the Poisson kernel. Therefore,
		\[
		S_{\varphi}(z,A_0) \equiv 1.
		\]
		Assume that $\Om$ is simply connected. Then,
		\[
		p(A_0,z) = \frac{\vert S(z,A_0)\vert^2}{S(A_0,A_0)}.
		\]
		Therefore, for $z\in\partial\Om$,
		\[
		L_{\varphi}(z,A_0) 
		=
		i \, \varphi(z) \, \overline{S_{\varphi}(z,A_0)} \, \overline{T(z)}
		=
		i \, \frac{\vert S(z,A_0)\vert^2}{S(A_0,A_0)} \, \overline{T(z)} = \frac{L(z,A_0) S(z,A_0)}{S(A_0,A_0)}.
		\]
		By the Identity principle,
		\[
		L_{\varphi}(z,A_0)  = \frac{L(z,A_0) \, S(z,A_0)}{S(A_0,A_0)}, \quad z\in\overline{\Om}.
		\]
		The weighted Ahlfors map $F_{\varphi}$ with respect to $A_0\in\Om$, is given by
		\[
		F_{\varphi}(z) = \frac{S(A_0,A_0)}{L(z,A_0) \, S(z,A_0)}.
		\]
		Note that the weighted Ahlfors map $F_{\varphi}\in \mathcal{B}_{\varphi}$ since $\Om$ is simply connected.
	\end{exam}
	
	\begin{exam}[see \cite{Nehari, zynda}]\label{mainexam}
		Let $\Omega\subset\mathbb{C}$ be a bounded $n$-connected domain with $C^{\infty}$ smooth boundary such that no boundary curve is an isolated point. Let 
		\[
		\varphi(z) = \vert f(z)\vert^2,
		\]
		where $f\in A^{\infty}(\Omega)$ is a non-vanishing function on $\overline{\Om}$.  Then
		\[
		S_{\varphi}(z,w) = \frac{1}{f(z)\overline{f(w)}} S(z,w)
		\quad\text{and}\quad
		L_{\varphi}(z,w)  = \frac{f(z)}{f(w)} L(z,w).
		\]
		Given $g\in H^2(\partial\Omega)$ and $a\in\Om$,
		\begin{eqnarray*}
			g(a) 
			&=&
			\frac{1}{f(a)} (gf)(a)
			=
			\frac{1}{f(a)} \int_{\partial\Om} (gf)(\zeta) \,\overline{S(\zeta,a)} \,ds
			=
			\frac{1}{f(a)} \int_{\partial\Om}(gf)(\zeta) \, \frac{\overline{S(\zeta,a)}}{\vert f(\zeta)\vert^2}\, \vert f(\zeta)\vert^2\, ds
			\\&=&
			\int_{\partial\Om}g(\zeta) \, \frac{\overline{S(\zeta,a)}}{\overline{f(\zeta)}f(a)}\, \varphi(\zeta)\, ds.
		\end{eqnarray*}
		Therefore, by the uniqueness of the weighted Szeg\H{o} kernel with respect to weight $\varphi$, we have
		\[
		S_{\varphi}(z,w) = \frac{1}{f(z)\overline{f(w)}} S(z,w),\quad z,w\in\Om.
		\]
		Now, for $z\in\partial\Om$ and $w\in\Om$,
		\[
		L_{\varphi}(z,w) 
		=
		i \,\vert f(z)\vert^2\, \overline{T(z)}\, \overline{S_{\varphi}(z,w)} 
		=
		i \,\vert f(z)\vert^2\, \overline{T(z)}\, \frac{1}{f(w)\overline{f(z)}}\,\overline{S(z,w)} 
		=
		\frac{f(z)}{f(w)} L(z,w).
		\]	
		The identity principle implies that
		\[
		L_{\varphi}(z,w)  = \frac{f(z)}{f(w)} L(z,w), 
		\quad z,w\in\Om.
		\]
		Thus, the weighted Szeg\H{o} kernel and the weighted Garabedian kernel inherit many nice properties of the classical Szeg\H{o} kernel and the Garabedian kernel, respectively.
		
		\medskip
		
		The weighted Ahlfors map $F_{\varphi}$ with respect to a point $a\in \Om$, is given by
		\[
		F_{\varphi}(z) = \frac{S_{\varphi}(z,a)}{L_{\varphi}(z,a)} 
		=
		\frac{f(a)^2}{\vert f(a)\vert^2} \, \frac{1}{f(z)^2} \frac{S(z,a)}{L(z,a)}.
		\]
		Thus, the weighted Ahlfors map $F_{\varphi}$ and the classical Ahlfors map $F$ with respect to a point $a\in\Om$, are related by
		\[
		F_{\varphi}(z) = \frac{f(a)^2}{\vert f(a)\vert^2} \, \frac{1}{f(z)^2} F(z).
		\]
		Here, $F_{\varphi}\in \mathcal{B}_{\varphi}$ for every $a\in\Om$, just as in the classical case.
		
		\medskip
		
		The function $F_{\varphi}$ maps the $n$ boundary curves $\gamma_i$ of $\Om$ onto $n$ Jordan closed curves $\Gamma_i$. Some of the images $\Gamma_i$ may coincide, for example all of them coincide if $\varphi\equiv 1$. The image of $\Om$ under $F_{\varphi}$ is open and connected and the boundary is the union of $\Gamma_i$. To see the image, $F_{\varphi}(\Om)$ is the connected component of 
		\[
		(\mathbb{C}\cup\{\infty\} )\setminus \cup_{i=1}^n\Gamma_i
		\]
		that contains $0$. The image is bounded as $F_{\varphi}$ has no poles.
	\end{exam}
	
	\begin{rem}
		Let $\Omega\subset\mathbb{C}$ be a bounded $n$-connected domain with $C^{\infty}$ smooth boundary such that no boundary curve is an isolated point. Let $\varphi$ be a positive real-valued $C^{\infty}$ smooth function on $\partial\Om$. It is known $($see \cite{UM, U-S}$)$ that the weighted Garabedian kernel $L_{\varphi}(z,w)$ does not vanish on $\overline{\Om}$ for all $w\in \Om$ if and only if there exists a non-vanishing function $f\in A^{\infty}(\Om)$ such that 
		\[
		L_{\varphi}(z,w) = L_{\vert f\vert^2}(z,w).
		\]
	\end{rem}
	
	\begin{exam}
		Let $\Omega\subset\mathbb{C}$ be a bounded $n$-connected domain with $C^{\infty}$ smooth boundary such that no boundary curve is an isolated point. For $A_0\in\Om$, let 
		\[
		\varphi(z) = \vert S(z,A_0)\vert^2.
		\]
		Given $g\in H^2(\partial\Omega)$,
		\begin{eqnarray*}
			g(A_0) 
			&=&
			\frac{1}{S(A_0,A_0)} \, g(A_0) \, S(A_0,A_0)
			=
			\frac{1}{S(A_0,A_0)} \int_{\partial\Om} g(\zeta) \, S(\zeta,A_0) \,\overline{S(\zeta,A_0)} \,ds
			\\&=&
			\frac{1}{S(A_0,A_0)} \int_{\partial\Om} g(\zeta) \,\varphi(\zeta) \,ds.
		\end{eqnarray*}
		Therefore, by the uniqueness of the weighted Szeg\H{o} kernel with respect to weight $\varphi$, we have
		\[
		S_{\varphi}(z,A_0) \equiv \frac{1}{S(A_0,A_0)}.
		\]
		Now, for $z\in\partial\Om$,
		\[
		L_{\varphi}(z,w) 
		=
		i \,\vert S(z,A_0)\vert^2\, \overline{T(z)}\, \overline{S_{\varphi}(z,w)} 
		=
		\frac{L(z,A_0)\,S(z,A_0)}{S(A_0,A_0)}.
		\]	
		The identity principle implies that
		\[
		L_{\varphi}(z,A_0)  = 	\frac{L(z,A_0)\,S(z,A_0)}{S(A_0,A_0)}, 
		\quad z\in\ov{\Om}.
		\]
		The weighted Ahlfors map $F_{\varphi}$ and the classical Ahlfors map $F$ with respect to $A_0\in \Om$, are related by
		\[
		F_{\varphi}(z) = \frac{1}{L(z,A_0)\,S(z,A_0)} = \frac{1}{S(z,A_0)^2} F(z).
		\]
        Here, $F_{\varphi}\in \mathcal{B}_{\varphi}$ if and only if $\Om$ is simply connected.
	\end{exam}
	
	\begin{exam}
		Let $\Omega\subset\mathbb{C}$ be a bounded $n$-connected domain with $C^{\infty}$ smooth boundary such that no boundary curve is an isolated point. For $A_0\in\Om$, let $B_0$ be a zero of $S(\cdot,A_0)$. Let 
		\[
		\varphi(z) = \frac{1}{\vert L(z,B_0)\vert^{2}}.
		\]
		Given $g\in H^2(\partial\Omega)$,
		\begin{eqnarray*}
			g(A_0) 
			&=&
			\frac{g(A_0)}{L(A_0,B_0)} \, L(A_0,B_0)
			=
			L(A_0,B_0) \int_{\partial\Om} \frac{g(\zeta)}{L(\zeta,B_0)} \,\overline{S(\zeta,A_0)} \,ds
			\\&=&
			L(A_0,B_0) \int_{\partial\Om}  \frac{g(\zeta)}{L(\zeta,B_0)} \,\overline{S(\zeta,A_0)} \,\vert L(\zeta,B_0)\vert^{2}\,\varphi(\zeta) \,ds
			\\&=&
			L(A_0,B_0) \int_{\partial\Om}  g(\zeta) \,\overline{S(\zeta,A_0)\,L(\zeta,B_0)} \,ds.
		\end{eqnarray*}
		Since $S(B_0,A_0) = 0$, the function $S(\cdot,A_0)\,L(\cdot,B_0) \in H^2(\partial\Om)$.
		Therefore, by the uniqueness of the weighted Szeg\H{o} kernel with respect to weight $\varphi$, we have
		\[
		S_{\varphi}(z,A_0) = S(z,A_0)\,L(z,B_0)\,\overline{L(A_0,B_0)}.
		\]
		Now, for $z\in\partial\Om$,
		\begin{eqnarray*}
			L_{\varphi}(z,A_0) 
			&=&
			i \, \frac{1}{\vert L(z,B_0)\vert^2} \, \overline{T(z)}\, \overline{S_{\varphi}(z,A_0)} 
			\\&=&
			i \, \frac{1}{\vert L(z,B_0)\vert^2} \, \overline{T(z)}\, \overline{S(z,A_0)}\,\overline{L(z,B_0)}\,L(A_0,B_0)
			\\&=&
			\frac{L(z,A_0)}{L(z,B_0)} \,L(A_0,B_0).
		\end{eqnarray*}
		The identity principle implies that
		\[
		L_{\varphi}(z,A_0)  = 	\frac{L(z,A_0)}{L(z,B_0)} \,L(A_0,B_0),
		\quad z\in\ov{\Om}.
		\]
		
		\medskip
		
		The weighted Ahlfors map $F_{\varphi}$ and the classical Ahlfors map $F$ with respect to $A_0\in \Om$, are related by
		\[
		F_{\varphi}(z) = \frac{\overline{L(A_0,B_0)}}{L(A_0,B_0)} \, L(z,B_0)^2 \, F(z).
		\]
		The weighted Ahlfors map $F_{\varphi}\in \mathcal{B}_{\varphi}$ if $S(z,A_0)$ vanishes at $z=B_0$ with multiplicity at least $2$. Otherwise, $F_{\varphi}$ has a simple pole at $z=B_0$.
	\end{exam}
	
	\begin{rem}
		\begin{enumerate}
			\item We do not know an example of a weight $\varphi$ for which $L_{\varphi}(z,a)$ or $S_{\varphi}(z,a)$ vanishes for some $z\in\partial\Om$ and $a\in\Om$. The problem is to find an example of such a weight $\varphi$ or prove that such a weight cannot exist.
			\item We do not know what happens when $\varphi(z) = \vert q(z)\vert^2$ where $q$ is an analytic function with at least one pole in $\Om$ extending $C^{\infty}$ smoothly to the boundary $\partial\Omega$.
		\end{enumerate}
	\end{rem}
	
	%%%%%%%%%%%%%%%%%%%%%%%%%%%%
	
	\section{Nehari's Construction of the weighted Szeg\H{o} and weighted Garabedian kernel}

As before, let $\Om\subset\mathbb{C}$ be a bounded $n$-connected domain with $C^{\infty}$ smooth boundary curves $\gamma_i$ and no isolated boundary point. Let $\varphi$ be a positive $C^{\infty}$ smooth function on $\partial\Om$. Theorem $1$ in \cite{Nehari} shows that there exists a meromorphic function $q$ on $\Omega$ that extends $C^{\infty}$-smoothly to $\pa \Omega$ such that the total number of zeros and poles counting multiplicities does not exceed $n-1$, and $\vert q(z) \vert^2 = \varphi(z)$ on $\pa \Om$. Some auxiliary functions arising in the proof of this theorem are of independent interest since they can be identified as the Szeg\H{o} kernel of certain closed subspaces of $H^2(\pa \Omega)$. To make all this clear, we give an exposition of Theorem $1$ in \cite{Nehari} for the sake of completeness and more importantly, to bring out the relevance and definition of the aforementioned auxiliary functions.

\medskip

For $a \in \Om$, recall that the Green's function $G(z, a)$ with pole at $z=a$ is defined as
\[
G(z,a) = -\log\vert z-a\vert + u_a(z)
\]
where $u_a(z)$ is the harmonic function of $z$ on $\Om$ that solves the Dirichlet problem with boundary data equal to $\log \vert z-a\vert$. Consider the harmonic function
\[
u(z) = -\frac{1}{4\pi} \int_{\partial\Om} \log \varphi(t) \frac{\partial G(z,t)}{\partial n}\, ds
\]
which has boundary values $2^{-1}\log \varphi(t)$. Let $2\pi p_{\nu}$, $\nu = 1,\ldots,n$, denote the periods of the harmonic conjugate $v$ of $u$. The function
\begin{equation}\label{Uform}
U(z) = u(z) + \sum_{\mu=1}^{n-1}\epsilon_{\mu} G(z,z_{\mu}),
\quad \quad z_{\mu}\in D, \mu =1,\ldots,n-1,
\end{equation}
where $\epsilon_{\mu}$ is either $+1$ or $-1$, has the same boundary values as $u(z)$. The period of $G(z,\z_{\mu})$ about $\gamma_{\nu}$ is
\[
\int_{\gamma_{\nu}} \frac{\partial G(z,z_{\mu})}{\partial n} ds 
=
(-2\pi) \frac{-1}{2\pi} \int_{\gamma_{\nu}} \frac{\partial G(z,z_{\mu})}{\partial n} ds
=
(-2\pi) \,\omega_{\nu}(z_{\mu})
\]
where $\omega_{\nu}$ denotes the harmonic function on $\Omega$ that takes the boundary values $\delta_{\nu\mu}$ on $\gamma_{\mu}$. The period $2\pi P_{\nu}$ of the harmonic conjugate $V$ of $U$ about $\gamma_{\nu}$ is given by
\[
P_{\nu} = p_{\nu} - \sum_{\mu=1}^{n-1} \epsilon_{\mu} \omega_{\nu}(z_{\mu}), 
\quad \quad \nu =1,\ldots,n.
\]
We want to choose $z_{\mu}$ and $\epsilon_{\mu}$ in such a way that $P_{\nu}$ are integers. Since $u$ is harmonic, $p_1+\ldots+p_n = 0$. Also, $\omega_1+\ldots+\omega_n \equiv 1$. So, 
\[
P_n = -\sum_{\nu=1}^{n-1} \left(p_{\nu} 
-
\sum_{\mu=1}^{n-1}\epsilon_{\mu}\omega_{\nu}(z_{\mu})\right) -\sum_{\mu=1}^{n-1}\epsilon_{\mu}
\]
and therefore, it is enough to make sure that $P_{\nu}$ are integers for $\nu=1,\ldots,n-1$. That is, we have to find $n-1$ points $z_1,\ldots,z_{n-1}$ in $\Om$ for which 
\begin{equation}\label{Ne1}
	\sum_{\mu=1}^{n-1}\epsilon_{\mu}\omega_{\nu}(z_{\mu}) = p_{\nu} + m_{\nu},
	\quad\quad \nu =1,\ldots, n-1
\end{equation}
holds for appropriately chosen integers $m_{\nu}$ and for a suitable choice of $\epsilon_{\mu}$. A schlicht (univalent) conformal mapping of $\Om$ transforms the harmonic measures of $\Om$ into the harmonic measures of its conformal image. Therefore, in order to solve (\ref{Ne1}), we can assume without loss of generality, that $\gamma_n$ is the real axis and $\Om$ is contained in the upper half plane. Since the harmonic functions $\omega_{\nu}$ vanish on the real axis, they can be analytically continued beyond the real axis by the identity $\omega_{\nu}(\bar{z}) = - \omega_{\nu}(z)$. The functions $\omega_{\nu}$ are thus harmonic in the entire domain $\Omega' = \Omega \cup \mathbb{R} \cup \tilde{\Om}$, where $\tilde{\Om}$ is the mirror image of $\Om$ with respect to the real axis. Since $\omega_{\nu}(\bar{z}) = - \omega_{\nu}(z)$, the system of equations (\ref{Ne1}) is equivalent to the system
\begin{equation}\label{Ne0}
	\sum_{\mu=1}^{n-1} \omega_{\nu}(z_{\mu}) = p_{\nu} + m_{\nu}, 
	\quad \quad \nu = 1,\ldots,n-1,
\end{equation}
if $z_{\mu}$ are not restricted to $\Om$ and can take any value in $\Omega'$. Now, consider the expression
\begin{equation}\label{Ne2}
	R = R(z_1,\ldots,z_{n-1})
	=
	\sum_{\nu =1}^{n-1}\left[\sum_{\mu=1}^{n-1} \omega_{\nu}(z_{\mu}) - p_{\nu} - m_{\nu}\right]^2.
\end{equation}
The function $R(z_1,\ldots,z_{n-1})$ is continuous with variables coming from $\Omega' + \Gamma'$ where $\Gamma'$ denotes the boundary of $\Omega'$. Thus, for a definite choice of integers $m_{\nu}$, $R$ has a non-negative minimum, which is attained for a set of $n-1$ points $z_{\mu}$ in $\Omega' + \Gamma'$. We now choose $m_{\nu}$ in such a way as to give this minimum its smallest possible value, which is possible because $-1\leq \omega_{\nu}(z)\leq 1$. If $m_1,\ldots,m_{n-1}$ are a choice of integers which give the smallest minimum of $R$, we write $p_{\nu} + m_{\nu} = \alpha_{\nu}$, and consider the solution of
\[
R = R(z_1,\ldots,z_{n-1}) 
=
\sum_{\nu =1}^{n-1}\left[\sum_{\mu=1}^{n-1} \omega_{\nu}(z_{\mu}) - \alpha_{\nu} \right]^2 = min,
\]
where $z_{\mu} \in \Omega'+\Gamma'$. Suppose $z_1\in \gamma_k$ ($k\neq n$). Since $\omega_k(z_1) = 1$, $\omega_{\nu}(z_1) = 0$ $(\nu\neq k)$, and $\omega_{\nu}(z)$ ($\nu =1,\ldots,n-1$) vanishes on the real axis, (\ref{Ne2}) will remain same if $z_1$ is transferred to the real axis and $m_k$ is replaced by $m_k-1$. Similarly, a point $z_{\mu}$ may be transferred from the mirror image of $\gamma_k$ to the real axis keeping (\ref{Ne2}) unchanged by replacing $m_k$ with $m_k + 1$. Therefore, there exists a minimizing set $S = \{z_{\mu}\}$ lying entirely in $\Om'$. Similarly, we can assume that $S$ does not contain both a point $z_{\mu}$ and its conjugate $\bar{z}_{\mu}$, because they can be replaced by points on the real axis without altering $(\ref{Ne2})$. Since $z_{\mu}$ minimize (\ref{Ne2}), we have
\[
0 =  \sum_{\nu =1}^{n-1}\left[\sum_{\mu=1}^{n-1} \omega_{\nu}(z_{\mu})-\alpha_{\nu}\right] {w'_{\nu}}(z_{\mu}) 
=
\sum_{\nu =1}^{n-1} a_{\nu} {w'_{\nu}}(z_{\mu}),
\quad \quad
a_{\nu} = \sum_{\mu=1}^{n-1} \omega_{\nu}(z_{\mu})-\alpha_{\nu}
\]
where $w_{\nu}(z)$ is (locally defined) holomorphic function whose real part is $\omega_{\nu}(z)$. Note that the derivative of $w_{\nu}$ is globally defined and equals $(\partial/\partial z) \omega_{\nu}$. Now, either all $a_{\nu}$ vanish or the derivative of the (locally defined) holomorphic function
\[
w(z) = \sum_{\nu=1}^{n-1} a_{\nu} w_{\nu}(z)
\] 
vanishes at $n-1$ points $z_1,\ldots,z_{n-1}$. In the first case, we have solved (\ref{Ne0}). So, assume that at least one of the constants $a_{\nu}$ is non-zero. The points of $S$ are critical points of $w(z)$. Since $a_{\nu}$ are real, we have ${w'_{\nu}}(\bar{z}) = - \overline{{w'_{\nu}}(z)}$. So, the set $S^*$ of points conjugate to those of $S$ is also a set of critical points of $w(z)$. If no point of $S$ is on the real axis, the sets $S$ and $S^*$ are disjoint and we thus have $2(n-1)$ critical points of $w(z)$, and $n-1$ of these points are in $\Om$. It is known that
\[
\overline{i \,\frac{\partial \omega_{\nu}}{\partial z}\, T(z)} = i \,\frac{\partial \omega_{\nu}}{\partial z} \,T(z)
\]
on $\partial\Omega$. Therefore, $i w'(z) T(z)$ is real on $\partial\Om$. By the argument principle, we have
\[
\text{no. of zeros of } w'(z) \text{ in } \Om = \frac{1}{2\pi} \Delta \arg w'(z) 
=
\frac{1}{2\pi} \left(\Delta\arg (i w'(z) T(z)) - \Delta\arg T(z)\right) = n-2,
\]
which is a contradiction. So, assume that one or more points of $S$ lie on the real axis. Note that $R$ is constant on the real axis. If $z_{\mu}$ is a point of $S$ lying on the real axis, then it can be replaced with any other point on the real axis without altering the value of the minimum. So, all the points of the real axis are critical points of $w(z)$. This means that
\[
\sum_{\nu=1}^{n-1} a_{\nu} {w'_{\nu}}(z) \equiv 0
\]
on the real axis, and by analytic continuation, throughout $\Om'$. But this is possible only if all $a_{\nu}$ vanish because ${w'_{\nu}}(z)$ are linearly independent. This leads to a contradiction and shows the solvability of (\ref{Ne0}).

\medskip

Therefore, the periods of the harmonic conjugate $V$ of $U$ are integral multiples of $2\pi$. Thus,
\[
q(z) = \exp\{U(z) + i V(z)\}
\]
is a single-valued meromorphic function on $\Om$ with simple zeros or simple poles at $z_{\mu}$, depending on whether $\epsilon_{\mu}$ is equal to $+1$ or $-1$. The number of zeros and poles of $q(z)$ counting multiplicity, does not exceed $n-1$ , and
\[
\vert q(z)\vert ^2 = \exp(2\,U(z)) = \exp\left(\log \varphi(z)\right) = \varphi(z),
\quad \quad z\in \partial\Omega.
\]
Note that if two of the points $z_{\mu}$ coincide but the respective $\epsilon_{\mu}$ are opposite in sign then the corresponding terms involving Green's function in (\ref{Uform}) get cancelled and the total number of zeros and poles of $q$ counting multiplicity will be less than $n-1$.

\medskip

Now, for $\zeta\in\Om$, consider the functions
\begin{equation}\label{Neq1}
	S^{(1)}(z,\zeta) = S(z,\zeta) + \sum_{\nu=1}^\sigma c_{\nu} S(z,z_{\nu}) + \sum_{\nu = \sigma +1}^{n-1} c_{\nu} L(z,z_{\nu}),
	\quad\quad \sigma\leq n-1,
\end{equation}
\begin{equation}\label{Neq2}
	L^{(1)}(z,\zeta) = L(z,\zeta) + \sum_{\nu=1}^\sigma \bar{c}_{\nu} L(z,z_{\nu}) + \sum_{\nu = \sigma +1}^{n-1} \bar{c}_{\nu} S(z,z_{\nu})
\end{equation}
where $z_1,\ldots,z_{\sigma}$ are the zeros of $q(z)$ and $z_{\sigma+1},\ldots,z_{n-1}$ are the poles of $q(z)$. The $n-1$ constants $c_{\nu}$ shall be determined by
\begin{equation}\label{Neq3}
	S^{(1)}(z_{\mu},\zeta) = 0, \quad \mu =1,\ldots, \sigma;
	\quad\quad
	L^{(1)}(z_{\mu}, \zeta) = 0, \quad \mu = \sigma+1,\ldots,n-1.
\end{equation}
This will be possible if the system of $n-1$ linear equations
\begin{multline}\label{Ne3}
	\sum_{\nu=1}^\sigma c_{\nu} S(z_{\mu},z_{\nu}) + \sum_{\nu = \sigma +1}^{n-1} c_{\nu} L(z_{\mu},z_{\nu})
	=
	-S(z_{\mu},\zeta),
	\quad\quad \mu = 1,\ldots,\sigma,
	\\
	\sum_{\nu=1}^\sigma c_{\nu} \overline{L(z_{\mu},z_{\nu})} + \sum_{\nu = \sigma +1}^{n-1} c_{\nu} \overline{S(z_{\mu},z_{\nu})} 
	=
	-\overline{L(z_{\mu},\zeta)},
	\quad\quad \mu = \sigma+1,\ldots,n-1,
\end{multline}
for the $n-1$ unknowns $c_1,\ldots,c_{n-1}$ has a solution. We may assume that not all the numbers $S(z_{\mu}, \zeta)$ $(\mu =1,\ldots,\sigma)$, $L(z_{\mu},\zeta)$ $(\mu = \sigma+1,\ldots,n-1)$ are zero. It is enough to show that the associated homogeneous system in $(\ref{Ne3})$ does not possess a non-trivial solution. On the contrary, if it had a non-trivial solution, there would exist constants $d_{\nu}$ $(\nu=1,\ldots, n)$, not all of which are zero, such that
\[
S^{(2)}(z_{\mu}) = 0,\quad \mu =1,\ldots, \sigma;\quad\quad
L^{(2)}(z_{\mu}) = 0, \quad \mu = \sigma+1,\ldots,n-1,
\]
where
\[
S^{(2)}(z) = \sum_{\nu =1}^{\sigma} d_{\nu} S(z,z_{\nu}) + \sum_{\nu = \sigma+1}^{n-1} d_{\nu} L(z,z_{\nu}),
\quad
L^{(2)}(z) = \sum_{\nu=1}^{\sigma} \bar{d}_{\nu} L(z,z_{\nu}) + \sum_{\nu=\sigma +1}^{n-1} \bar{d}_{\nu} S(z,z_{\nu}).
\]
If two of the points $z_{\nu}$ coincide, at the point $z_{\tau}$, $\tau\leq \sigma$, then (\ref{Neq1}) and (\ref{Neq2}) will each have an extra term involving $\partial S(z,z_{\tau})/\partial \bar{z}_{\tau}$ and $\partial L(z,z_{\tau})/\partial z_{\tau}$, respectively. The condition in (\ref{Neq3}) which is lost by the two points coinciding is replaced by the condition $\partial S^{(1)}(z_{\tau}, \zeta)/\partial z_{\tau} = 0$. If $\tau > \sigma$, then the roles of the functions $S$ and $L$ are interchanged in modifying (\ref{Neq1}), (\ref{Neq2}) and the condition (\ref{Neq3}). Coincidence of more than two points will result in the appearance of higher derivatives in the appropriate places. These changes do not affect the validity of the considerations which follow. We have
\[
\overline{S^{(2)}(z)} = \frac{1}{i} L^{(2)}(z) T(z),
\quad \quad z\in \partial \Om.
\]
Therefore,
\[
\frac{1}{i} S^{(2)}(z)L^{(2)}(z) T(z) \geq 0,
\quad\quad z\in\partial\Omega,
\]
which implies that
\[
\frac{1}{i} \int_{\partial\Om}S^{(2)}(z)L^{(2)}(z) dz >0. 
\]
By Cauchy's theorem, the function $S^{(2)}(z)L^{(2)}(z)$ cannot be analytic in $\Om$, which is a contradiction. Therefore, we have proved that ($\ref{Ne3}$) has a unique solution. Furthermore,
\begin{equation}\label{S1L1}
\overline{S^{(1)}(z,\zeta)} = \frac{1}{i} L^{(1)}(z,\zeta) T(z),
\quad\quad z\in \partial\Om.
\end{equation}
The weighted Szeg\H{o} and Garabedian kernels are given by
\begin{equation}
	S_{\varphi}(z,\zeta) = \frac{S^{(1)}(z,\zeta)}{q(z)\, \overline{q(\zeta)}},
	\quad\quad
	L_{\varphi}(z,\zeta) = \frac{L^{(1)}(z,\zeta) \,q(z)}{q(\zeta)}.
\end{equation}
If $q(z)$ happens to have a zero or a pole at $z=\zeta$, the above definitions of $S_{\varphi}(z,\zeta)$ and $L_{\varphi}(z,\zeta)$ need to be slightly modified. If $\eta$ is a point close to $\zeta$, we define
\begin{equation}
S_{\varphi}(z,\zeta) = \lim_{\eta\rightarrow \zeta} S_{\varphi}(z,\eta),
\quad\quad 
L_{\varphi}(z,\zeta) = \lim_{\eta\rightarrow\zeta} L_{\varphi}(z,\eta).
\end{equation}
This completes the exposition of Nehari's theorem and yields expressions for the weighted Szeg\H{o} and Garabedian kernels in terms of the auxiliary functions $S^{(1)}(z, \zeta), L^{(1)}(z, \zeta)$. 

\subsubsection{\bf{Observation 1}}\label{SC1}

If $\Omega$ is simply connected $(n=1)$, then $q$ does not have any zero or pole in $\Omega$. That is $q\in A^{\infty}(\Om)$ and is non-vanishing. Therefore, it follows from Example \ref{mainexam} that
\begin{equation}
S_{\varphi}(z,w) = \frac{S(z,w)}{q(z)\,\overline{q(w)}},
\quad\quad \text{and} \quad\quad
L_{\varphi}(z,w) = \frac{L(z,w) \,q(z)}{q(w)}.
\end{equation}
This recovers a result already stated in \cite{Nehari, UM} and gives an explicit expression for the weighted Szeg\H{o} kernel and the weighted Garabedian kernel for any bounded simply connected domain $\Om$, given that we have the Riemann map $F:\Om\rightarrow\mathbb{D}$ and the function $q$ above.

\subsubsection{\bf{Observation 2}}

\medskip

Let $\Om$ be a bounded $n$-connected domain with $C^{\infty}$ smooth boundary and no isolated boundary point. Let $\varphi$ be a positive $C^{\infty}$ smooth function on $\partial\Om$ and $q$ as in Nehari's theorem with $\vert q \vert^2 = \varphi$ on $\pa \Om$.

\medskip

Assume that $q$ does not have any pole in $\Om$, that is, $q\in A^{\infty}(\Om)$. Let $z_1,\ldots,z_{\sigma}$ be the zeros of $q$ with multiplicity $m_1,\ldots,m_{\sigma}$. We have $m_1 + \ldots+m_{\sigma}\leq n-1$. Consider the \textit{effective divisor} $A=\{(z_1,m_1),\ldots,(z_{\sigma},m_{\sigma})\}$ and define 
\begin{equation}
H_{A}^2(\Omega) = \{f\in H^2(\partial\Omega): f(z_{\mu}) = 0 \text{ with multiplicity }m_{\mu}, \,\,\mu = 1,\ldots,\sigma\}
\end{equation}
The space $H_{A}^2(\Omega)$ is a closed subspace of $H^2(\partial\Om)$ and therefore it is a Hilbert space. For $w\in\Om$, the evaluation functionals 
\[
E_w: H_{A}^2(\Om) \ni f \mapsto f(w) \in \mathbb{C}
\]
are continuous and therefore by Riesz Representation theorem, there exists a unique function $S_{A}(\cdot,w)\in H_{A}^2(\Om)$ such that
\begin{equation}
f(w) = \int_{z\in\partial\Om} f(z) \,\overline{S_{A}(z,w)}\, ds.
\end{equation}
The function $S_{A}(z,w)$ is the reproducing kernel of $H_{A}^2(\Om)$. For $f\in H^2_{\varphi}(\partial\Om)$, the function $f.q \in H_{A}(\Om)$. Therefore, for $w\in\Om$,
\[
f(w) q(w) = \int_{z\in\partial\Om}  f(z)\,q(z) \,\overline{S_{A}(z,w)}\,ds 
=
\int_{z\in\partial\Om}  f(z)\,\frac{\overline{S_{A}(z,w)}}{\overline{q(z)}}\,\varphi(z)\,ds.
\]
If $q(w)\neq 0$, it follows by the uniqueness of the weighted Szeg\H{o} kernel that
\begin{equation}
S_{\varphi}(z,w) = \frac{S_{A}(z,w)}{q(z)\,\overline{q(w)}}.
\end{equation}
In the case $q(w) = 0$, if $\eta$ is a point close to $w$, we have
\begin{equation}
S_{\varphi}(z,w) = \lim\limits_{\eta\rightarrow w} S_{\varphi}(z,\eta) = 
 \lim\limits_{\eta\rightarrow w} \frac{S_{A}(z,\eta)}{q(z)\,\overline{q(\eta)}}.
\end{equation}
Thus, the function $S^{(1)}(z,w)$ in Nehari's construction is the function $S_{A}(z,w)$, when $q$ is assumed to have no poles. 

\medskip

Now, let $v\in H_A^2(\Om)^{\perp}$. For $u\in L^2(\partial\Om)$, the function $q\,\mathcal{C}u \in H_A^2(\Omega)$. Therefore,
\[
0=\langle q\,\mathcal{C}u, v\rangle
= \langle \mathcal{C}u, v\bar{q}\rangle
= \langle u, \mathcal{C}^* (v\bar{q})\rangle
\]
This implies that $ \mathcal{C}^* (v\bar{q}) = 0$, that is, $v\bar{q} = \ov{T}\, \overline{\mathcal{C}(\bar{v}q\ov{T})}$. Thus, $v \bar{q} = \overline{TH}$ for some $H\in H^2(\partial\Om)$. Conversely, let $v\in L^2(\partial\Om)$ be such that $v \bar{q} = \overline{TH}$ for some $H\in H^2(\partial\Om)$. Then, for any $u\in H_A^2(\Om)$
\[
\langle u, v \rangle = \left\langle \frac{u}{q}, v\bar{q} \right\rangle 
=  \left\langle \frac{u}{q}, \overline{TH}\right\rangle = 0
\]
because $u/q\in H^2(\partial\Omega)$. Thus, every $v\in H^2_A(\Om)^{\perp}$ satisfies $v\bar{q} = \overline{TH}$ for some $H\in H^2(\partial\Om)$. Let $P_A$ denote the orthogonal projection of $L^2(\partial\Om)$ onto $H^2(\partial\Om)$. The orthogonal decomposition of $u\in L^2(\partial\Om)$ satisfies
\[
u\, \bar{q} = (P_A u) \,\bar{q} + \overline{TH}
\]
where $H = P_A(\bar{u}q\overline{T})$ because $\bar{u}\,q\,\overline{T} = H + \overline{(P_A u)}\, q\, \overline{T}$ is the orthogonal decomposition of $\bar{u}q\overline{T}$.

\medskip

Now, the orthogonal projection of the Cauchy kernel $C_a$ onto $H_A^2(\Om)$ is $S_A(\cdot,a)$ because
\[
P_A C_a (\zeta) = \int_{z\in\partial \Om} C_a(z) S_A(\zeta,z) ds = S_A(\zeta,a)
\]
by the Cauchy integral formula. Therefore, writing the orthogonal decomposition of $C_a$ gives
\[
\overline{\left(\frac{1}{2\pi i} \frac{T(\zeta)}{\zeta - a}\right)}\, \ov{q(\zeta)}
=
C_a(\zeta)\, \ov{q(\zeta)}
= S_A(\zeta,a)\, \ov{q(\zeta)} + \overline{T(\zeta) H_a(\zeta)}
\]
where $H_a\in H^2(\partial\Om)$ and $H_a = P_A(\ov{C_a} q \ov{T})$. We now define
\begin{equation}
L_A(\zeta,a) = i \left( \frac{1}{2\pi i} \frac{1}{\zeta - a} - \frac{H_a(\zeta)}{q(\zeta)}\right).
\end{equation}
By the definition, it follows that for $\zeta\in\partial\Om$,
\begin{eqnarray*}
\overline{S_A(\zeta,a)} = \left( \frac{1}{2\pi i} \frac{1}{\zeta - a} - \frac{H_a(\zeta)}{q(\zeta)} \right) T(\zeta) = \frac{1}{i} L_A(\zeta,a) T(\zeta).
\end{eqnarray*}
That is,
\begin{equation}
\overline{S_A(\zeta,a)} = \frac{1}{i} L_A(\zeta,a) T(\zeta),
\quad \quad \zeta\in\partial\Om, a\in\Om.
\end{equation}
A comparison with (\ref{S1L1}) and analytic continuation implies that the function $L^{(1)}(z,w)$ in Nehari's construction is the function $L_A(z,w)$, when $q$ is assumed to have no poles. Also,
\begin{equation}
L_{\varphi}(z,w) = \frac{L_A(z,w) q(z)}{q(w)}
\end{equation}
when $q(w)\neq 0$. In the case $q(w) = 0$, if $\eta$ is a point close to $w$, we have
\begin{equation}
L_{\varphi}(z,w) = \lim\limits_{\eta\rightarrow w} L_{\varphi}(z,\eta) 
=
\lim\limits_{\eta\rightarrow w}\frac{L_A(z,\eta) q(z)}{q(\eta)}.
\end{equation}
The above arguments are true even if $q$ is assumed to be any function such that $\vert q(z)\vert^2 = \varphi(z)$ on $\partial\Om$, without any restriction on the zeros.
To summarize,

\begin{thm}\label{neh}
Let $\Om$ be a bounded $n$-connected domain with $C^{\infty}$ smooth boundary and no isolated boundary point. Let $\varphi$ be a positive $C^{\infty}$ smooth function on $\partial\Om$ and $q\in A^{\infty}(\Om)$ be any function such that
\[
\vert q(z)\vert^2 = \varphi(z),
\quad\quad z\in\partial\Om.
\]
Let $z_1,\ldots,z_{\sigma}$ be the zeros of $q$ with multiplicity $m_1,\ldots,m_{\sigma}$. Let $A=\{(z_1,m_1),\ldots,(z_{\sigma},m_{\sigma})\}$. Corresponding to the set $A$, let $S_{A}(z,w)$ denote the reproducing kernel of the Hilbert space
\[
H_{A}^2(\Omega) = \{f\in H^2(\partial\Omega): f(z_{\mu}) = 0 \text{ with multiplicity }m_{\mu}, \,\,\mu = 1,\ldots,\sigma\}.
\]
Then, the weighted Szeg\H{o} kernel $S_{\varphi}$ is given by
\[
S_{\varphi}(z,w) = \frac{S_A(z,w)}{q(z)\ov{q(w)}}
\quad \text{ if }w\neq z_\mu,
\quad\quad
S_{\varphi}(z,z_{\mu}) = \lim\limits_{\eta\rightarrow z_{\mu}}  S_{\varphi}(z,\eta).
\]
The orthogonal projection of the Cauchy kernel $C_a$ onto $H_A^2(\Om)$ is $S_A(\cdot,a)$. And, the orthogonal decomposition of $C_a$ satisfies
\[
C_a \bar{q}= S_A(\cdot, a) \bar{q} + \overline{TH_a}, 
\quad\quad \text{for some }H_a\in H^2(\partial\Om).
\]
We define the function $L_A$ by
\[
L_A(\zeta,a) = i \left(\frac{1}{2\pi i}\frac{1}{\zeta - a}- \frac{H_a(\zeta)}{q(\zeta)}\right).
\]
The functions $S_A$ and $L_A$ are related by
\[
\overline{S_A(\zeta,a)} = \frac{1}{i} L_A(\zeta,a) T(\zeta),
\quad \quad \zeta\in\partial\Om, a\in\Om.
\]
The weighted Garabedian kernel $L_{\varphi}$ is given by
\[
L_{\varphi}(z,w) = \frac{L_A(z,w) q(z)}{q(w)}
\quad \text{ if } w\neq z_{\mu},
\quad\quad
L_{\varphi}(z,z_{\mu}) = \lim\limits_{\eta\rightarrow z_{\mu}} L_{\varphi}(z,\eta).
\]
\end{thm}

\begin{rem} 
	It is worthwhile looking at Theorem 6 and Theorem 7 in \cite{Nehari}. In Theorem 6, the extremal function $F$ on domain $D$ that maximizes $\vert f'(\zeta)\vert$ for $\zeta\in D$ where $f\in\mathcal{O}(D)$, $\vert f(z)\vert\leq 1$ for $z\in D$ and $f(\zeta) = f(\alpha_1) = \cdots=f(\alpha_m)=0$ $(\alpha_1,\ldots,\alpha_m\in D)$, is same as
	\begin{equation}
	F(z) = \frac{S_A(z,\zeta)}{L_A(z,\zeta)}
	\end{equation}
	for $A = \{\alpha_1,\ldots,\alpha_m\}$. This can be directly checked using the definitions of $S_A$ and $L_A$.
\end{rem}

\begin{cor}\label{SC2}
	Let $\Om$ be a bounded simply connected domain with $C^{\infty}$ smooth boundary. Let $A=\{(z_1,m_1),\ldots,(z_{\sigma},m_{\sigma})\}$ where $z_i\in\Om$ and $m_i$ are positive integers. Let $p\in A^{\infty}(\Om)$ be a function such that $p(z_i) = 0$ with multiplicity $m_i$ for $i=1,\ldots,\sigma$, and $p$ does not vanish elsewhere. Let 
	\[
	\varphi(z) = \vert p(z)\vert^2, \quad z\in\partial\Om.
	\]
	There exists a nonvanishing function $q\in A^{\infty}(\Om)$ such that $\vert q(z)\vert^2 = \varphi(z)$ for $z\in\partial\Om$. Note that the function $q$ is unique modulo a multiplicative constant of modulus $1$. Then,
	\begin{equation}
	S_A(z,w) = \frac{p(z) \overline{p(w)}}{q(z)\overline{q(w)}} \,S(z,w)
	\quad
	\text{and}
	\quad
	L_A(z,w) = \frac{q(z)p(w)}{p(z)q(w)} \,L(z,w).
	\end{equation}
\end{cor}
\begin{proof}
	It follows using Observation $1$ in \ref{SC1} and Theorem \ref{neh}. If $q_1$ and $q_2$ are two candidates for $q$ then $q_1/q_2$ is a non-vanishing holomorphic function from $\Om$ into $\mathbb{D}$ with modulus $1$ on $\partial\Om$. This forces $q_1/q_2$ to be a constant.
\end{proof}

\begin{exam}
	As a very simple example for Corollary \ref{SC2}, let $\Om = \mathbb{D}$ and $A=\{(0, n)\}$ where $n$ is a positive integer. Then $p(z) = z^n$ and $q(z) \equiv 1$. We thus have
	\[
	S_A(z,w) = z^n \bar{w}^n S_{\mathbb{D}}(z,w) = \frac{1}{2\pi} \frac{z^n\bar{w}^n}{1-z\bar{w}}
	\quad
	\text{and}
	\quad
	L_{A}(z,w) = \frac{w^n}{z^n} L_{\mathbb{D}}(z,w) = \frac{1}{2\pi} \frac{w^n}{z^n} \frac{1}{z-w}.
	\]
\end{exam}


%%%%%%%%%%%%%%%%%%%%%%%%%%%%%%%%%%%%%%%%%%%
	
	
	\section{Weighted Carath\'eodory metric and properties}
	
	The Carathéodory metric $c_{\Om}$ on a smoothly bounded domain $\Om\subset\mathbb{C}$ is defined by
	\begin{equation}
	c_{\Om}(z)=\sup\{\vert f'(z)\vert:\, f:\Om\rightarrow\mathbb{D}\,\text{ is holomorphic and }\,f(z)=0\},\quad z\in \Om.
	\end{equation}
	For $z\in \Om$, the solutions of the above extremal problem are of the form $e^{i\theta} f_z$, where $f_z$ denotes the Ahlfors map of $\Om$ based at point $z$. It is known that
	\[
	f_z(w)=\frac{S(w,z)}{L(w,z)},\quad w\in \Om.
	\]
	Since $L(\cdot,z)$ has a pole at $w=z$ with residue $1/2\pi$, we obtain
	\[
	{f_z}'(z)=\lim_{w\rightarrow z}\frac{f_z(w)-f_z(z)}{w-z}=\lim_{w\rightarrow z}\frac{f_z(w)}{w-z}=\lim_{w\rightarrow z}\frac{S(w,z)}{L(w,z) (w-z)}=\frac{\lim_{w\rightarrow z} S(w,z)}{1/2\pi}=2\pi S(z,z).
	\]
	Therefore, 
	\begin{equation}
	c_{\Om}(z)=2\pi\, S(z,z).
	\end{equation}
	
	\begin{defn}[Weighted Carathéodory metric]
		Let $\Omega\subset\mathbb{C}$ be a bounded $n$-connected domain with $C^{\infty}$ smooth boundary such that no boundary curve is an isolated point. Let $\varphi$ be a positive real-valued $C^{\infty}$ smooth function on $\partial\Om$. We define the weighted Carathéodory metric $c_{\varphi}$ with respect to weight $\varphi$ by
		\begin{equation}
		c_{\varphi}(z) = 2\pi\, S_{\varphi}(z,z)
		\end{equation}
		for $z\in\Om$.
	\end{defn}
	
	Let there exist a function $q\in A^{\infty}(\Om)$ such that $\vert q(z)\vert^2 = \varphi(z)$ for $z\in\partial\Om$. Let $z_1,\ldots,z_{\sigma}$ be the zeros of $q$ with multiplicity $m_1,\ldots,m_{\sigma}$. Let $A=\{(z_1,m_1),\ldots,(z_{\sigma},m_{\sigma})\}$. Then,
	\begin{equation}
	c_{\varphi}(z) = 2\pi S_{\varphi}(z,z) = 2\pi \frac{S_A(z,z)}{\vert q(z)\vert^2},\quad\quad z\neq z_{\mu}.
	\end{equation}

    %%%%%%%%%%%%%%%%%%%%
	\section{Weighted Szeg\H{o} kernels in terms of the classical kernels}
	
	Let $\Omega\subset\mathbb{C}$ be a bounded $n$-connected domain with $C^{\infty}$ smooth boundary such that no boundary curve is an isolated point. Let $\gamma_j$, $j=1,\ldots,n$ denote the $n$ boundary curves of $\Omega$. The harmonic measure functions $\omega_j$ are unique harmonic functions on $\Omega$ that take value $1$ on $\gamma_j$ and $0$ on $\gamma_i$ for $i\neq j$. Let $F_j'$ denote the holomorphic functions on $\Omega$ given by $(1/2)(\partial/\partial z)\omega_j(z)$.
	
	\medskip
	
	Let $\varphi$ be a positive real-valued $C^{\infty}$ smooth function on $\partial\Om$. It is known (see \cite{Nehari_1}) that 
	\begin{equation}\label{1}
	S_{\varphi}(z,w) S_{1/\varphi}(z,w) = S(z,w)^2 + \sum_{\nu=1}^{n-1} a_{\nu}(w)F_{\nu}'(z)
	\end{equation}
	\begin{equation}\label{2}
	L_{\varphi}(z,w) L_{1/\varphi}(z,w) = L(z,w)^2 + \sum_{\nu=1}^{n-1} \overline{a_{\nu}(w)}F_{\nu}'(z) 
	\end{equation}
	where $a_{\nu}$ are constants depending upon $w\in\Om$.
	
	\medskip
	
	It is known (see \cite{Be0, Bergman}) that if $K$ denotes the classical Bergman kernel of $\Om$ then
	\begin{equation}\label{ks}
	K(z,w) = 4\pi S(z,w)^2 + \sum_{j,k=1}^{n-1} c_{jk}F_{j}'(z) \overline{F_{k}'(w)}
	\end{equation}
	for some constants $c_{jk}$. Infact, if $a\in\Om$ is such that $S(\cdot,a)$ has $n-1$ distinct zeros $a_1,\ldots,a_{n-1}$ then (see \cite{Be0, Be5})
	\begin{equation}
	K(z,a) = 4\pi S(z,a)^2 + 2\pi\sum_{j=1}^{n-1} \frac{K(a_j,a)}{S'(a_j,a)}L(z,a_j)S(z,a).
	\end{equation}
	This is because the $n-1$ functions $L(z,a_j)S(z,a)$ form a basis for $\mathcal{F}'=span_{\mathbb{C}}\{ F_j'(z) : j=1,\ldots, n-1 \}$ (a theorem due to Schiffer, see \cite{Be0, Be5, Sch}). This result due to Schiffer also implies that
	\begin{equation}\label{ks2}
	K(z,w) = 4\pi S(z,w)^2 + \sum_{j,k=1}^{n-1} \lambda_{jk} L(z,a_j)S(z,a) \overline{L(w,a_k)S(w,a)}
	\end{equation}
	for some constants $\lambda_{jk}$ (see \cite{Be4} as well).
	
	\medskip
	
Coming back to (\ref{1}), we will see that
\begin{equation}\label{aid}
a_{\nu}(w) = \sum_{\mu=1}^{n-1} A_{\mu\nu}\overline{F_{\mu}'(w)}
\end{equation}
for some constants $A_{\mu\nu}$. This can be seen following an argument in \cite[see page 137]{Be0}. For $k=1,\ldots,n-1$
\[
\int_{\gamma_{k}} (S_{\varphi}(z,w) S_{1/\varphi}(z,w) - S(z,w)^2 ) \,dz
=
\sum_{\nu=1}^{n-1}a_{\nu}(w)\int_{\gamma_{k}} F_{\nu}'(z)\,dz 
=
\sum_{\nu=1}^{n-1} P_{k\nu} \,a_{\nu}(w)
\]
where $P=[P_{\mu\nu}]$ for $\mu,\nu=1,\ldots,n-1$, denotes the nonsingular matrix of periods 
\[
\int_{\gamma_{\mu}} F_{\nu}'(z) \,dz.
\]
Let $\tilde{\gamma}_{k}$ be a curve that is homotopic to $\gamma_{k}$, but that is inside $\Om$. Then
\begin{equation}\label{3}
\int_{\tilde{\gamma}_{k}} (S_{\varphi}(z,w) S_{1/\varphi}(z,w) - S(z,w)^2 ) \,dz
=
\sum_{\nu=1}^{n-1} P_{k\nu} \,a_{\nu}(w).
\end{equation}
Note that $a_{\nu}(w)$ are antiholomorphic functions in $w$. Now,
\begin{eqnarray*}
S_{\varphi}(z,w) S_{1/\varphi}(z,w) - S(z,w)^2 
&=& 
\overline{S_{\varphi}(w,z) S_{1/\varphi}(w,z) - S(w,z)^2}
= \overline{\sum_{\nu=1}^{n-1} a_{\nu}(z) F_{\nu}'(w)} 
\\&=& \sum_{\nu=1}^{n-1} \overline{a_{\nu}(z)} \overline{F_{\nu}'(w)}.
\end{eqnarray*}
The integral in (\ref{3}) can be approximated by a finite Riemann sum which will therefore lie in the linear span of functions $\{\overline{F_{\nu}'(w)}\}_{\nu=1}^{n-1}$. A limiting argument will show that the integral also lies in the linear span of $\overline{F_{\nu}'(w)}$. Thus, we have
\[
\sum_{\nu=1}^{n-1} Q_{k\nu} \overline{F_{\nu}'(w)} = 
\sum_{\nu=1}^{n-1} P_{k\nu} \,a_{\nu}(w),
\quad \quad k=1,\ldots,n-1
\]
for some constants $Q_{k\nu}$. Let $Q$ denote the matrix $[Q_{\mu\nu}]$. Therefore,
\[
P \,\cdot\,[a_{\nu}(w)]_{\nu=1}^{n-1}  = Q \,\cdot\,[\overline{F_{\nu}'(w)}]_{\nu=1}^{n-1}.
\]
Since $P$ is nonsingular, we obtain (\ref{aid}).
Following the same line of reasoning for (\ref{2}), we have
\begin{equation}\label{Neq4}
	S_{\varphi}(z,w) S_{1/\varphi}(z,w) = S(z,w)^2 + \sum_{j,k=1}^{n-1} a_{jk}\,F_{j}'(z)\overline{F_k'(w)}
\end{equation}
\begin{equation}\label{Neq5}
	L_{\varphi}(z,w) L_{1/\varphi}(z,w) = L(z,w)^2 + \sum_{j,k=1}^{n-1} \overline{a_{jk}}\, F_{j}'(z) F_{k}'(w) 
\end{equation}
for some constants $a_{jk}$. This also implies that
\begin{equation}\label{Neq6}
	S_{\varphi}(z,w) S_{1/\varphi}(z,w) = S(z,w)^2 + \sum_{j,k=1}^{n-1} \alpha_{jk}\,L(z,a_j)S(z,a) \overline{L(w,a_k)S(w,a)}
\end{equation}
\begin{equation}\label{Neq7}
	L_{\varphi}(z,w) L_{1/\varphi}(z,w) = L(z,w)^2 + \sum_{j,k=1}^{n-1} \overline{\alpha_{jk}}\, L(z,a_j)S(z,a) L(w,a_k)S(w,a)
\end{equation}
for some constants $\alpha_{jk}$.
Relations (\ref{Neq1}), (\ref{Neq2}), Theorem \ref{neh}, relations (\ref{Neq4}) -- (\ref{Neq7}) together may help to shed more light on the properties of the weighted Szeg\H{o} and weighted Garabedian kernel.

\medskip

Finally, we also observe from the relations (\ref{ks}), (\ref{ks2}), (\ref{Neq4}) and (\ref{Neq6}) that
\begin{equation}
4\pi\,S_{\varphi}(z,w) S_{1/\varphi}(z,w) = K(z,w) + \sum_{j,k=1}^{n-1} b_{jk}\,F_{j}'(z)\overline{F_k'(w)}
\end{equation}
\begin{equation}
4\pi\, S_{\varphi}(z,w) S_{1/\varphi}(z,w) = K(z,w) + \sum_{j,k=1}^{n-1} \beta_{jk}\,L(z,a_j)S(z,a) \overline{L(w,a_k)S(w,a)}
\end{equation}
for some constants $b_{jk}$ and $\beta_{jk}$.

\medskip

Thus, the weighted Szeg\H{o} kernels corresponding to weights $\varphi$ and $1/\varphi$ can be expressed in terms of the classical kernel functions.
    
	
	\begin{thebibliography}{1}
		
		\bibitem{Be0}
		Bell, Steven R.: \textit{The Cauchy transform, potential theory and conformal mapping}. Second edition. Chapman and Hall/CRC, Boca Raton, FL, 2016.
		
		\bibitem{BeKf} 
		Bell Steven R.; Kaleem, Faisal:
		\textit{The structure of the semigroup of proper holomorphic mappings of a planar domain to the unit disc.}
		Comput. Methods Funct. Theory \textbf{8} (2008), no. 1-2, 225–242. 
		
		\bibitem{Be1}
		Bell, Steven R.: \textit{Complexity in complex analysis.}
		Adv. Math. \textbf{172} (2002), no. 1, 15–52.
		
		% \bibitem{Be2}
		% Bell, Steven R.: \textit{The fundamental role of the Szegő kernel in potential theory and complex analysis.}
		% J. Reine Angew. Math. \textbf{525} (2000), 1–16. 
		
		% \bibitem{Be3}
		% Bell, Steven R.: \textit{Ahlfors maps, the double of a domain, and complexity in potential theory and conformal mapping.}
		% J. Anal. Math. \textbf{78} (1999), 329–344. 
		
		\bibitem{Be4}
		Bell, Steven R.:  \textit{Complexity of the classical kernel functions of potential theory.}
		Indiana Univ. Math. J. \textbf{44} (1995), no. 4, 1337–1369.
		
		\bibitem{Be5} 
		Bell, Steven R.: \textit{The Szegő projection and the classical objects of potential theory in the plane}.
		Duke Math. J. \textbf{64} (1991), no. 1, 1–26.
		
		\bibitem{Bergman}
		Bergman, Stefan: \textit{The kernel function and conformal mapping.}
		Second, revised edition. Mathematical Surveys, No. \textbf{V}. American Mathematical Society, Providence, R.I., 1970.
		
%		\bibitem{Burbea}
%		Burbea, Jacob: \textit{The higher order curvatures of weighted span metrics on Riemann surfaces.}
%		Arch. Math. (Basel) \textbf{43} (1984), no. 5, 473–479. 

	 %    \bibitem{G2}
	 %    Grunsky, Helmut:
	 %    \textit{Lectures on theory of functions in multiply connected domains.}
	 %    Vandenhoeck and Ruprecht, Göttingen, 1978, 253 pp.

		% \bibitem{G1}
		% Grunsky, Helmut: \textit{Über die konforme Abbildung mehrfach zusammenhängender Bereiche auf mehrblättrige Kreise. II.} Abh. Preuss. Akad. Wiss. Math.-Nat. Kl. 1941 (1941), no. \textbf{11}, 8 pp. 
		
		\bibitem{JV}
		Jain, Aakanksha; Verma, Kaushal: \textit{Weighted Szegő kernels on planar domains}. Illinois J. Math. \textbf{69} (2025), no. 2, 235–268.
		
		\bibitem{Nehari_1}
		Nehari, Zeev:
		\textit{On weighted kernels}. J. Analyse Math. \textbf{2} (1952), 126–149. 
		
		\bibitem{Nehari}
		Nehari, Zeev: \textit{A class of domain functions and some allied extremal problems.} 
		Trans. Amer. Math. Soc. \textbf{69} (1950), 161–178.
		
		\bibitem{pw-um}
		Pasternak-Winiarski, Zbigniew; Zynda, Tomasz Lukasz:
		\textit{Weighted Szegő kernels}. Geometric methods in physics \textbf{XXXV}, 145–151, Trends Math., Birkhäuser/Springer, Cham, 2018. 
		
%		\bibitem{ra1}
%		Pasternak-Winiarski, Zbigniew; Wójcicki, Paweł: \textit{Weighted generalization of the Ramadanov's theorem and further considerations.}
%		Czechoslovak Math. J. \textbf{68}(143) (2018), no. 3, 829–842. 
%		
%		\bibitem{pw1}
%		Pasternak-Winiarski, Zbigniew: \textit{On weights which admit the reproducing kernel of Bergman type.}
%		Internat. J. Math. Math. Sci. \textbf{15} (1992), no. 1, 1–14.
%		
%		\bibitem{pw2} 
%		Pasternak-Winiarski, Zbigniew: \textit{On the dependence of the reproducing kernel on the weight of integration.}
%		J. Funct. Anal. \textbf{94} (1990), no. 1, 110–134. 

		\bibitem{Sch} Schiffer, Menahem:
		\textit{Various types of orthogonalization}. Duke Math. J. \textbf{17} (1950), 329–366.
		
		\bibitem{UM}
		Uehara, Masahiro: \textit{The Nehari problem for the weighted Szegő kernels}. Reproducing kernels and their applications (Newark, DE, 1997), 213–221, Int. Soc. Anal. Appl. Comput., \textbf{3}, Kluwer Acad. Publ., Dordrecht, 1999. 
		
		\bibitem{U-S}
		Uehara, Masahiro; Saitoh, Saburou:
	\textit{Some remarks for the weighted Szegő kernel functions}. Math. Japon. \textbf{29} (1984), no. 6, 887–891.
		
		\bibitem{zynda}
		Żynda, Tomasz Łukasz:
		\textit{Weighted generalization of the Szegö kernel and how it can be used to prove general theorems of complex analysis}. Geometric methods in physics \textbf{XXXVII}, 212–218, Trends Math., Birkhäuser/Springer, Cham, 2019.
		
		
	\end{thebibliography}
	
\end{document}
